% Options for packages loaded elsewhere
% Options for packages loaded elsewhere
\PassOptionsToPackage{unicode}{hyperref}
\PassOptionsToPackage{hyphens}{url}
\PassOptionsToPackage{dvipsnames,svgnames,x11names}{xcolor}
%
\documentclass[
  croatian,
  letterpaper,
  DIV=11,
  numbers=noendperiod]{scrartcl}
\usepackage{xcolor}
\usepackage{amsmath,amssymb}
\setcounter{secnumdepth}{5}
\usepackage{iftex}
\ifPDFTeX
  \usepackage[T1]{fontenc}
  \usepackage[utf8]{inputenc}
  \usepackage{textcomp} % provide euro and other symbols
\else % if luatex or xetex
  \usepackage{unicode-math} % this also loads fontspec
  \defaultfontfeatures{Scale=MatchLowercase}
  \defaultfontfeatures[\rmfamily]{Ligatures=TeX,Scale=1}
\fi
\usepackage{lmodern}
\ifPDFTeX\else
  % xetex/luatex font selection
\fi
% Use upquote if available, for straight quotes in verbatim environments
\IfFileExists{upquote.sty}{\usepackage{upquote}}{}
\IfFileExists{microtype.sty}{% use microtype if available
  \usepackage[]{microtype}
  \UseMicrotypeSet[protrusion]{basicmath} % disable protrusion for tt fonts
}{}
\makeatletter
\@ifundefined{KOMAClassName}{% if non-KOMA class
  \IfFileExists{parskip.sty}{%
    \usepackage{parskip}
  }{% else
    \setlength{\parindent}{0pt}
    \setlength{\parskip}{6pt plus 2pt minus 1pt}}
}{% if KOMA class
  \KOMAoptions{parskip=half}}
\makeatother
% Make \paragraph and \subparagraph free-standing
\makeatletter
\ifx\paragraph\undefined\else
  \let\oldparagraph\paragraph
  \renewcommand{\paragraph}{
    \@ifstar
      \xxxParagraphStar
      \xxxParagraphNoStar
  }
  \newcommand{\xxxParagraphStar}[1]{\oldparagraph*{#1}\mbox{}}
  \newcommand{\xxxParagraphNoStar}[1]{\oldparagraph{#1}\mbox{}}
\fi
\ifx\subparagraph\undefined\else
  \let\oldsubparagraph\subparagraph
  \renewcommand{\subparagraph}{
    \@ifstar
      \xxxSubParagraphStar
      \xxxSubParagraphNoStar
  }
  \newcommand{\xxxSubParagraphStar}[1]{\oldsubparagraph*{#1}\mbox{}}
  \newcommand{\xxxSubParagraphNoStar}[1]{\oldsubparagraph{#1}\mbox{}}
\fi
\makeatother

\usepackage{color}
\usepackage{fancyvrb}
\newcommand{\VerbBar}{|}
\newcommand{\VERB}{\Verb[commandchars=\\\{\}]}
\DefineVerbatimEnvironment{Highlighting}{Verbatim}{commandchars=\\\{\}}
% Add ',fontsize=\small' for more characters per line
\usepackage{framed}
\definecolor{shadecolor}{RGB}{241,243,245}
\newenvironment{Shaded}{\begin{snugshade}}{\end{snugshade}}
\newcommand{\AlertTok}[1]{\textcolor[rgb]{0.68,0.00,0.00}{#1}}
\newcommand{\AnnotationTok}[1]{\textcolor[rgb]{0.37,0.37,0.37}{#1}}
\newcommand{\AttributeTok}[1]{\textcolor[rgb]{0.40,0.45,0.13}{#1}}
\newcommand{\BaseNTok}[1]{\textcolor[rgb]{0.68,0.00,0.00}{#1}}
\newcommand{\BuiltInTok}[1]{\textcolor[rgb]{0.00,0.23,0.31}{#1}}
\newcommand{\CharTok}[1]{\textcolor[rgb]{0.13,0.47,0.30}{#1}}
\newcommand{\CommentTok}[1]{\textcolor[rgb]{0.37,0.37,0.37}{#1}}
\newcommand{\CommentVarTok}[1]{\textcolor[rgb]{0.37,0.37,0.37}{\textit{#1}}}
\newcommand{\ConstantTok}[1]{\textcolor[rgb]{0.56,0.35,0.01}{#1}}
\newcommand{\ControlFlowTok}[1]{\textcolor[rgb]{0.00,0.23,0.31}{\textbf{#1}}}
\newcommand{\DataTypeTok}[1]{\textcolor[rgb]{0.68,0.00,0.00}{#1}}
\newcommand{\DecValTok}[1]{\textcolor[rgb]{0.68,0.00,0.00}{#1}}
\newcommand{\DocumentationTok}[1]{\textcolor[rgb]{0.37,0.37,0.37}{\textit{#1}}}
\newcommand{\ErrorTok}[1]{\textcolor[rgb]{0.68,0.00,0.00}{#1}}
\newcommand{\ExtensionTok}[1]{\textcolor[rgb]{0.00,0.23,0.31}{#1}}
\newcommand{\FloatTok}[1]{\textcolor[rgb]{0.68,0.00,0.00}{#1}}
\newcommand{\FunctionTok}[1]{\textcolor[rgb]{0.28,0.35,0.67}{#1}}
\newcommand{\ImportTok}[1]{\textcolor[rgb]{0.00,0.46,0.62}{#1}}
\newcommand{\InformationTok}[1]{\textcolor[rgb]{0.37,0.37,0.37}{#1}}
\newcommand{\KeywordTok}[1]{\textcolor[rgb]{0.00,0.23,0.31}{\textbf{#1}}}
\newcommand{\NormalTok}[1]{\textcolor[rgb]{0.00,0.23,0.31}{#1}}
\newcommand{\OperatorTok}[1]{\textcolor[rgb]{0.37,0.37,0.37}{#1}}
\newcommand{\OtherTok}[1]{\textcolor[rgb]{0.00,0.23,0.31}{#1}}
\newcommand{\PreprocessorTok}[1]{\textcolor[rgb]{0.68,0.00,0.00}{#1}}
\newcommand{\RegionMarkerTok}[1]{\textcolor[rgb]{0.00,0.23,0.31}{#1}}
\newcommand{\SpecialCharTok}[1]{\textcolor[rgb]{0.37,0.37,0.37}{#1}}
\newcommand{\SpecialStringTok}[1]{\textcolor[rgb]{0.13,0.47,0.30}{#1}}
\newcommand{\StringTok}[1]{\textcolor[rgb]{0.13,0.47,0.30}{#1}}
\newcommand{\VariableTok}[1]{\textcolor[rgb]{0.07,0.07,0.07}{#1}}
\newcommand{\VerbatimStringTok}[1]{\textcolor[rgb]{0.13,0.47,0.30}{#1}}
\newcommand{\WarningTok}[1]{\textcolor[rgb]{0.37,0.37,0.37}{\textit{#1}}}

\usepackage{longtable,booktabs,array}
\usepackage{calc} % for calculating minipage widths
% Correct order of tables after \paragraph or \subparagraph
\usepackage{etoolbox}
\makeatletter
\patchcmd\longtable{\par}{\if@noskipsec\mbox{}\fi\par}{}{}
\makeatother
% Allow footnotes in longtable head/foot
\IfFileExists{footnotehyper.sty}{\usepackage{footnotehyper}}{\usepackage{footnote}}
\makesavenoteenv{longtable}
\usepackage{graphicx}
\makeatletter
\newsavebox\pandoc@box
\newcommand*\pandocbounded[1]{% scales image to fit in text height/width
  \sbox\pandoc@box{#1}%
  \Gscale@div\@tempa{\textheight}{\dimexpr\ht\pandoc@box+\dp\pandoc@box\relax}%
  \Gscale@div\@tempb{\linewidth}{\wd\pandoc@box}%
  \ifdim\@tempb\p@<\@tempa\p@\let\@tempa\@tempb\fi% select the smaller of both
  \ifdim\@tempa\p@<\p@\scalebox{\@tempa}{\usebox\pandoc@box}%
  \else\usebox{\pandoc@box}%
  \fi%
}
% Set default figure placement to htbp
\def\fps@figure{htbp}
\makeatother



\ifLuaTeX
\usepackage[bidi=basic]{babel}
\else
\usepackage[bidi=default]{babel}
\fi
% get rid of language-specific shorthands (see #6817):
\let\LanguageShortHands\languageshorthands
\def\languageshorthands#1{}


\setlength{\emergencystretch}{3em} % prevent overfull lines

\providecommand{\tightlist}{%
  \setlength{\itemsep}{0pt}\setlength{\parskip}{0pt}}



 


\KOMAoption{captions}{tableheading}
\makeatletter
\@ifpackageloaded{tcolorbox}{}{\usepackage[skins,breakable]{tcolorbox}}
\@ifpackageloaded{fontawesome5}{}{\usepackage{fontawesome5}}
\definecolor{quarto-callout-color}{HTML}{909090}
\definecolor{quarto-callout-note-color}{HTML}{0758E5}
\definecolor{quarto-callout-important-color}{HTML}{CC1914}
\definecolor{quarto-callout-warning-color}{HTML}{EB9113}
\definecolor{quarto-callout-tip-color}{HTML}{00A047}
\definecolor{quarto-callout-caution-color}{HTML}{FC5300}
\definecolor{quarto-callout-color-frame}{HTML}{acacac}
\definecolor{quarto-callout-note-color-frame}{HTML}{4582ec}
\definecolor{quarto-callout-important-color-frame}{HTML}{d9534f}
\definecolor{quarto-callout-warning-color-frame}{HTML}{f0ad4e}
\definecolor{quarto-callout-tip-color-frame}{HTML}{02b875}
\definecolor{quarto-callout-caution-color-frame}{HTML}{fd7e14}
\makeatother
\makeatletter
\@ifpackageloaded{caption}{}{\usepackage{caption}}
\AtBeginDocument{%
\ifdefined\contentsname
  \renewcommand*\contentsname{Table of contents}
\else
  \newcommand\contentsname{Table of contents}
\fi
\ifdefined\listfigurename
  \renewcommand*\listfigurename{List of Figures}
\else
  \newcommand\listfigurename{List of Figures}
\fi
\ifdefined\listtablename
  \renewcommand*\listtablename{List of Tables}
\else
  \newcommand\listtablename{List of Tables}
\fi
\ifdefined\figurename
  \renewcommand*\figurename{Figure}
\else
  \newcommand\figurename{Figure}
\fi
\ifdefined\tablename
  \renewcommand*\tablename{Table}
\else
  \newcommand\tablename{Table}
\fi
}
\@ifpackageloaded{float}{}{\usepackage{float}}
\floatstyle{ruled}
\@ifundefined{c@chapter}{\newfloat{codelisting}{h}{lop}}{\newfloat{codelisting}{h}{lop}[chapter]}
\floatname{codelisting}{Listing}
\newcommand*\listoflistings{\listof{codelisting}{List of Listings}}
\makeatother
\makeatletter
\makeatother
\makeatletter
\@ifpackageloaded{caption}{}{\usepackage{caption}}
\@ifpackageloaded{subcaption}{}{\usepackage{subcaption}}
\makeatother
\usepackage{bookmark}
\IfFileExists{xurl.sty}{\usepackage{xurl}}{} % add URL line breaks if available
\urlstyle{same}
\hypersetup{
  pdftitle={Tjedan 6: Programiranje u R-u},
  pdflang={hr},
  colorlinks=true,
  linkcolor={blue},
  filecolor={Maroon},
  citecolor={Blue},
  urlcolor={Blue},
  pdfcreator={LaTeX via pandoc}}


\title{Tjedan 6: Programiranje u R-u}
\usepackage{etoolbox}
\makeatletter
\providecommand{\subtitle}[1]{% add subtitle to \maketitle
  \apptocmd{\@title}{\par {\large #1 \par}}{}{}
}
\makeatother
\subtitle{Funkcije, uvjeti i ponovljive analize}
\author{}
\date{2025-03-29}
\begin{document}
\maketitle

\renewcommand*\contentsname{Table of contents}
{
\hypersetup{linkcolor=}
\setcounter{tocdepth}{3}
\tableofcontents
}

\begin{Shaded}
\begin{Highlighting}[]
\FunctionTok{library}\NormalTok{(tidyverse)}
\end{Highlighting}
\end{Shaded}

\begin{tcolorbox}[enhanced jigsaw, coltitle=black, arc=.35mm, colframe=quarto-callout-note-color-frame, left=2mm, bottomrule=.15mm, toprule=.15mm, leftrule=.75mm, bottomtitle=1mm, colback=white, title=\textcolor{quarto-callout-note-color}{\faInfo}\hspace{0.5em}{Ishodi učenja}, opacityback=0, colbacktitle=quarto-callout-note-color!10!white, toptitle=1mm, titlerule=0mm, opacitybacktitle=0.6, breakable, rightrule=.15mm]

Nakon ovog predavanja moći ćete

\begin{enumerate}
\def\labelenumi{\arabic{enumi}.}
\tightlist
\item
  Objasniti zašto su vlastite funkcije korisne za izbjegavanje
  ponavljanja koda i smanjenje grešaka.
\item
  Napisati vlastitu R funkciju s argumentima i podrazumijevanim
  (default) vrijednostima.
\item
  Koristiti uvjetne naredbe (\texttt{if}, \texttt{else},
  \texttt{if\_else()}, \texttt{case\_when()}) za kontrolu toka programa.
\item
  Koristiti \texttt{for} petlje za ponavljanje operacija nad skupom
  elemenata.
\item
  Koristiti \texttt{map()} funkcije iz paketa purrr kao modernu
  alternativu petljama.
\item
  Primijeniti principe DRY (Don't Repeat Yourself) na pisanje
  analitičkih skripti.
\item
  Organizirati analitičku skriptu s jasnom strukturom: učitavanje,
  čišćenje, analiza, vizualizacija, izvoz.
\item
  Prepoznati kada je pisanje vlastite funkcije isplativije od kopiranja
  koda.
\end{enumerate}

\end{tcolorbox}

\section{Koliko programiranja treba
komunikolog?}\label{koliko-programiranja-treba-komunikolog}

Ovo pitanje zaslužuje iskren odgovor. Ne trebate postati softverski
inženjer. Ne trebate znati pisati web aplikacije, baze podataka ili
algoritme strojnog učenja. Ali trebate znati dovoljno programiranja da
vaše analize budu \textbf{ponovljive}, \textbf{prilagodljive} i
\textbf{manje podložne greškama}.

Zamislite sljedeću situaciju. Radite analizu medijskih navika za
klijenta. Napravili ste čišćenje podataka, deskriptivnu statistiku, osam
grafova i izvještaj. Klijent je zadovoljan, ali tjedan dana kasnije
kaže: ``Dobili smo još 200 odgovora na anketu, možete li ponoviti
analizu s novim podacima?'' Ako ste sve radili ručno u Excelu, to znači
ponoviti svaki korak od nule. Ako ste napisali R skriptu, to znači
promijeniti jednu liniju koda (putanju do nove datoteke) i pokrenuti
skriptu. Pet sekundi umjesto pet sati.

Programiranje u kontekstu analize podataka nije apstraktno akademsko
znanje. To je praktična vještina koja vas čini bržima, preciznijima i
profesionalnijima. U ovom tjednu naučit ćemo tri temeljne programerske
koncepte (funkcije, uvjetne naredbe i iteraciju) i pokazati kako ih
koristiti u kontekstu koji je relevantan za komunikologe.

\begin{center}\rule{0.5\linewidth}{0.5pt}\end{center}

\section{Naši podaci: newsletter kampanje}\label{sec-podaci}

Ovaj tjedan koristimo dataset o 50 newsletter kampanja jednog
informativnog portala. Za svaku kampanju imamo podatke o tipu, stilu
naslova, vremenu slanja, broju pretplatnika, open rateu (postotak
otvaranja), click rateu (postotak klikova) i drugim metrikama.

\begin{Shaded}
\begin{Highlighting}[]
\NormalTok{nl }\OtherTok{\textless{}{-}} \FunctionTok{read\_csv}\NormalTok{(}\StringTok{"../resources/datasets/newsletter\_campaigns.csv"}\NormalTok{)}
\FunctionTok{glimpse}\NormalTok{(nl)}
\end{Highlighting}
\end{Shaded}

\begin{verbatim}
Rows: 50
Columns: 13
$ campaign_id      <chr> "NL-001", "NL-002", "NL-003", "NL-004", "NL-005", "NL~
$ campaign_type    <chr> "special_report", "weekly_digest", "special_report", ~
$ subject_style    <chr> "personalizirani", "hitno", "upitni", "informativni",~
$ day_sent         <chr> "petak", "petak", "utorak", "ponedjeljak", "utorak", ~
$ send_hour        <dbl> 8, 11, 8, 9, 20, 16, 13, 6, 11, 19, 11, 19, 9, 15, 18~
$ subscribers      <dbl> 11770, 14266, 10652, 23113, 9847, 9150, 23450, 12798,~
$ open_rate        <dbl> 0.2410, 0.2696, 0.3023, 0.2134, 0.2887, 0.1921, 0.288~
$ click_rate       <dbl> 0.0858, 0.0065, 0.0309, 0.0656, 0.0273, 0.0651, 0.015~
$ unsubscribe_rate <dbl> 0.00447, 0.00283, 0.00519, 0.00000, 0.00082, 0.00391,~
$ word_count       <dbl> 499, 378, 437, 545, 559, 146, 428, 210, 309, 519, 376~
$ n_links          <dbl> 2, 9, 4, 1, 3, 1, 10, 3, 8, 4, 9, 3, 8, 7, 5, 7, 1, 5~
$ has_image        <lgl> TRUE, TRUE, FALSE, TRUE, TRUE, TRUE, TRUE, TRUE, FALS~
$ revenue          <dbl> 0.00, 0.00, 0.00, 0.00, 0.00, 0.00, 0.00, 4659.53, 0.~
\end{verbatim}

\begin{Shaded}
\begin{Highlighting}[]
\NormalTok{nl }\SpecialCharTok{|\textgreater{}}
  \FunctionTok{count}\NormalTok{(campaign\_type, }\AttributeTok{sort =} \ConstantTok{TRUE}\NormalTok{)}
\end{Highlighting}
\end{Shaded}

\begin{verbatim}
# A tibble: 5 x 2
  campaign_type      n
  <chr>          <int>
1 special_report    17
2 weekly_digest     12
3 sponsored          8
4 breaking_news      7
5 event_promo        6
\end{verbatim}

Ovo je manji dataset od prethodnih tjedana, ali upravo to ga čini
pogodnim za učenje programiranja. S 50 redova možete vidjeti svaki korak
i razumjeti što vaš kod radi.

\begin{center}\rule{0.5\linewidth}{0.5pt}\end{center}

\section{Zašto funkcije? Problem kopiranja
koda}\label{sec-zasto-funkcije}

Krenimo od konkretnog problema. Recimo da za svaki tip kampanje želite
izračunati sažetak s prosjekom, medijanom i standardnom devijacijom open
ratea. Jedan pristup je kopiranje koda.

\begin{Shaded}
\begin{Highlighting}[]
\CommentTok{\# Sažetak za weekly digest}
\NormalTok{nl }\SpecialCharTok{|\textgreater{}}
  \FunctionTok{filter}\NormalTok{(campaign\_type }\SpecialCharTok{==} \StringTok{"weekly\_digest"}\NormalTok{) }\SpecialCharTok{|\textgreater{}}
  \FunctionTok{summarise}\NormalTok{(}
    \AttributeTok{n =} \FunctionTok{n}\NormalTok{(),}
    \AttributeTok{or\_prosjek =} \FunctionTok{round}\NormalTok{(}\FunctionTok{mean}\NormalTok{(open\_rate), }\DecValTok{3}\NormalTok{),}
    \AttributeTok{or\_medijan =} \FunctionTok{round}\NormalTok{(}\FunctionTok{median}\NormalTok{(open\_rate), }\DecValTok{3}\NormalTok{),}
    \AttributeTok{or\_sd =} \FunctionTok{round}\NormalTok{(}\FunctionTok{sd}\NormalTok{(open\_rate), }\DecValTok{3}\NormalTok{)}
\NormalTok{  )}
\end{Highlighting}
\end{Shaded}

\begin{verbatim}
# A tibble: 1 x 4
      n or_prosjek or_medijan or_sd
  <int>      <dbl>      <dbl> <dbl>
1    12      0.289       0.29 0.066
\end{verbatim}

\begin{Shaded}
\begin{Highlighting}[]
\CommentTok{\# Sažetak za breaking news (kopiran kod s jednom promjenom)}
\NormalTok{nl }\SpecialCharTok{|\textgreater{}}
  \FunctionTok{filter}\NormalTok{(campaign\_type }\SpecialCharTok{==} \StringTok{"breaking\_news"}\NormalTok{) }\SpecialCharTok{|\textgreater{}}
  \FunctionTok{summarise}\NormalTok{(}
    \AttributeTok{n =} \FunctionTok{n}\NormalTok{(),}
    \AttributeTok{or\_prosjek =} \FunctionTok{round}\NormalTok{(}\FunctionTok{mean}\NormalTok{(open\_rate), }\DecValTok{3}\NormalTok{),}
    \AttributeTok{or\_medijan =} \FunctionTok{round}\NormalTok{(}\FunctionTok{median}\NormalTok{(open\_rate), }\DecValTok{3}\NormalTok{),}
    \AttributeTok{or\_sd =} \FunctionTok{round}\NormalTok{(}\FunctionTok{sd}\NormalTok{(open\_rate), }\DecValTok{3}\NormalTok{)}
\NormalTok{  )}
\end{Highlighting}
\end{Shaded}

\begin{verbatim}
# A tibble: 1 x 4
      n or_prosjek or_medijan or_sd
  <int>      <dbl>      <dbl> <dbl>
1     7      0.198      0.195 0.055
\end{verbatim}

Ovo radi, ali ima tri problema. Prvo, ako želite promijeniti izračun
(recimo dodati trimmed mean), morate to napraviti na svakom mjestu gdje
ste kopirali kod. Drugo, svako kopiranje je prilika za grešku. Možda
zaboravite promijeniti ime kampanje na jednom mjestu. Treće, kad imate
pet ili deset tipova kampanja, kod postaje nepregledano dugačak.

Naravno, za ovaj specifični problem znamo elegantno rješenje:
\texttt{group\_by()}.

\begin{Shaded}
\begin{Highlighting}[]
\NormalTok{nl }\SpecialCharTok{|\textgreater{}}
  \FunctionTok{group\_by}\NormalTok{(campaign\_type) }\SpecialCharTok{|\textgreater{}}
  \FunctionTok{summarise}\NormalTok{(}
    \AttributeTok{n =} \FunctionTok{n}\NormalTok{(),}
    \AttributeTok{or\_prosjek =} \FunctionTok{round}\NormalTok{(}\FunctionTok{mean}\NormalTok{(open\_rate), }\DecValTok{3}\NormalTok{),}
    \AttributeTok{or\_medijan =} \FunctionTok{round}\NormalTok{(}\FunctionTok{median}\NormalTok{(open\_rate), }\DecValTok{3}\NormalTok{),}
    \AttributeTok{or\_sd =} \FunctionTok{round}\NormalTok{(}\FunctionTok{sd}\NormalTok{(open\_rate), }\DecValTok{3}\NormalTok{),}
    \AttributeTok{.groups =} \StringTok{"drop"}
\NormalTok{  )}
\end{Highlighting}
\end{Shaded}

\begin{verbatim}
# A tibble: 5 x 5
  campaign_type      n or_prosjek or_medijan or_sd
  <chr>          <int>      <dbl>      <dbl> <dbl>
1 breaking_news      7      0.198      0.195 0.055
2 event_promo        6      0.25       0.23  0.086
3 special_report    17      0.259      0.255 0.067
4 sponsored          8      0.248      0.246 0.043
5 weekly_digest     12      0.289      0.29  0.066
\end{verbatim}

Ali \texttt{group\_by()} ne rješava svaki problem. Kad trebate ponoviti
složeniju analizu (koja uključuje čišćenje, više izračuna, graf i
tablicu) za različite podskupove podataka, vlastite funkcije postaju
nezamjenjive.

\begin{center}\rule{0.5\linewidth}{0.5pt}\end{center}

\section{Pisanje vlastite funkcije}\label{sec-funkcije}

Funkcija u R-u je objekt koji prima ulazne podatke (argumente), izvršava
niz operacija i vraća rezultat. Već koristite funkcije svaki dan:
\texttt{mean()}, \texttt{filter()}, \texttt{ggplot()} su sve funkcije
koje je netko napisao. Sad ćete naučiti pisati vlastite.

\subsection{Anatomija funkcije}\label{anatomija-funkcije}

\begin{Shaded}
\begin{Highlighting}[]
\CommentTok{\# Funkcija koja pretvara postotke u razlomke}
\NormalTok{postotak\_u\_razlomak }\OtherTok{\textless{}{-}} \ControlFlowTok{function}\NormalTok{(postotak) \{}
\NormalTok{  postotak }\SpecialCharTok{/} \DecValTok{100}
\NormalTok{\}}

\FunctionTok{postotak\_u\_razlomak}\NormalTok{(}\DecValTok{25}\NormalTok{)}
\end{Highlighting}
\end{Shaded}

\begin{verbatim}
[1] 0.25
\end{verbatim}

\begin{Shaded}
\begin{Highlighting}[]
\FunctionTok{postotak\_u\_razlomak}\NormalTok{(}\FloatTok{73.5}\NormalTok{)}
\end{Highlighting}
\end{Shaded}

\begin{verbatim}
[1] 0.735
\end{verbatim}

Raščlanimo sintaksu. \texttt{postotak\_u\_razlomak} je ime funkcije (kao
ime bilo kojeg objekta, dodjeljujemo ga s \texttt{\textless{}-}).
Ključna riječ \texttt{function()} govori R-u da kreiramo funkciju.
Unutar zagrada su argumenti (ulazni podatci). Unutar vitičastih zagrada
\texttt{\{\}} je tijelo funkcije (operacije koje se izvršavaju). Zadnji
izraz u tijelu je povratna vrijednost (ono što funkcija vraća).

\subsection{Funkcija s više
argumenata}\label{funkcija-s-viux161e-argumenata}

\begin{Shaded}
\begin{Highlighting}[]
\CommentTok{\# Funkcija za izračun engagement ratea}
\NormalTok{engagement\_rate }\OtherTok{\textless{}{-}} \ControlFlowTok{function}\NormalTok{(clicks, opens) \{}
\NormalTok{  rate }\OtherTok{\textless{}{-}}\NormalTok{ clicks }\SpecialCharTok{/}\NormalTok{ opens}
  \FunctionTok{round}\NormalTok{(rate, }\DecValTok{4}\NormalTok{)}
\NormalTok{\}}

\FunctionTok{engagement\_rate}\NormalTok{(}\AttributeTok{clicks =} \DecValTok{150}\NormalTok{, }\AttributeTok{opens =} \DecValTok{1200}\NormalTok{)}
\end{Highlighting}
\end{Shaded}

\begin{verbatim}
[1] 0.125
\end{verbatim}

\begin{Shaded}
\begin{Highlighting}[]
\FunctionTok{engagement\_rate}\NormalTok{(}\AttributeTok{clicks =} \DecValTok{80}\NormalTok{, }\AttributeTok{opens =} \DecValTok{500}\NormalTok{)}
\end{Highlighting}
\end{Shaded}

\begin{verbatim}
[1] 0.16
\end{verbatim}

Funkcija prima dva argumenta i vraća zaokruženi omjer. Kad pozivate
funkciju, argumente možete navesti po imenu (\texttt{clicks\ =\ 150})
ili po poziciji. Po imenu je sigurnije jer nije bitno kojim redoslijedom
ih navedete.

\subsection{Default vrijednosti
argumenata}\label{default-vrijednosti-argumenata}

Ponekad želite da argument ima podrazumijevanu (default) vrijednost koju
korisnik može promijeniti ako želi.

\begin{Shaded}
\begin{Highlighting}[]
\CommentTok{\# Funkcija za sažetak numeričke varijable}
\NormalTok{sazetak\_varijable }\OtherTok{\textless{}{-}} \ControlFlowTok{function}\NormalTok{(x, }\AttributeTok{decimale =} \DecValTok{2}\NormalTok{) \{}
  \FunctionTok{tibble}\NormalTok{(}
    \AttributeTok{n =} \FunctionTok{length}\NormalTok{(x),}
    \AttributeTok{n\_NA =} \FunctionTok{sum}\NormalTok{(}\FunctionTok{is.na}\NormalTok{(x)),}
    \AttributeTok{prosjek =} \FunctionTok{round}\NormalTok{(}\FunctionTok{mean}\NormalTok{(x, }\AttributeTok{na.rm =} \ConstantTok{TRUE}\NormalTok{), decimale),}
    \AttributeTok{medijan =} \FunctionTok{round}\NormalTok{(}\FunctionTok{median}\NormalTok{(x, }\AttributeTok{na.rm =} \ConstantTok{TRUE}\NormalTok{), decimale),}
    \AttributeTok{sd =} \FunctionTok{round}\NormalTok{(}\FunctionTok{sd}\NormalTok{(x, }\AttributeTok{na.rm =} \ConstantTok{TRUE}\NormalTok{), decimale),}
    \AttributeTok{min =} \FunctionTok{round}\NormalTok{(}\FunctionTok{min}\NormalTok{(x, }\AttributeTok{na.rm =} \ConstantTok{TRUE}\NormalTok{), decimale),}
    \AttributeTok{max =} \FunctionTok{round}\NormalTok{(}\FunctionTok{max}\NormalTok{(x, }\AttributeTok{na.rm =} \ConstantTok{TRUE}\NormalTok{), decimale)}
\NormalTok{  )}
\NormalTok{\}}

\CommentTok{\# Korištenje s default decimala (2)}
\FunctionTok{sazetak\_varijable}\NormalTok{(nl}\SpecialCharTok{$}\NormalTok{open\_rate)}
\end{Highlighting}
\end{Shaded}

\begin{verbatim}
# A tibble: 1 x 7
      n  n_NA prosjek medijan    sd   min   max
  <int> <int>   <dbl>   <dbl> <dbl> <dbl> <dbl>
1    50     0    0.25    0.25  0.07  0.11  0.41
\end{verbatim}

\begin{Shaded}
\begin{Highlighting}[]
\CommentTok{\# Korištenje s 4 decimale}
\FunctionTok{sazetak\_varijable}\NormalTok{(nl}\SpecialCharTok{$}\NormalTok{open\_rate, }\AttributeTok{decimale =} \DecValTok{4}\NormalTok{)}
\end{Highlighting}
\end{Shaded}

\begin{verbatim}
# A tibble: 1 x 7
      n  n_NA prosjek medijan     sd   min   max
  <int> <int>   <dbl>   <dbl>  <dbl> <dbl> <dbl>
1    50     0   0.255   0.252 0.0676 0.115 0.408
\end{verbatim}

Argument \texttt{decimale\ =\ 2} ima default vrijednost 2. Ako ga ne
navedete pri pozivu, koristi se 2. Ako ga eksplicitno navedete, koristi
se vaša vrijednost. Ovo čini funkciju fleksibilnom bez opterećivanja
korisnika nepotrebnim odlukama.

\subsection{Funkcija koja radi s
tibbleom}\label{funkcija-koja-radi-s-tibbleom}

Funkcije koje primaju cijeli tibble i koriste dplyr glagole unutar sebe
su izuzetno korisne u praksi.

\begin{Shaded}
\begin{Highlighting}[]
\CommentTok{\# Funkcija za sažetak kampanje po tipu}
\NormalTok{sazetak\_kampanje }\OtherTok{\textless{}{-}} \ControlFlowTok{function}\NormalTok{(data, tip) \{}
\NormalTok{  data }\SpecialCharTok{|\textgreater{}}
    \FunctionTok{filter}\NormalTok{(campaign\_type }\SpecialCharTok{==}\NormalTok{ tip) }\SpecialCharTok{|\textgreater{}}
    \FunctionTok{summarise}\NormalTok{(}
      \AttributeTok{tip =}\NormalTok{ tip,}
      \AttributeTok{n =} \FunctionTok{n}\NormalTok{(),}
      \AttributeTok{prosjek\_or =} \FunctionTok{round}\NormalTok{(}\FunctionTok{mean}\NormalTok{(open\_rate), }\DecValTok{3}\NormalTok{),}
      \AttributeTok{prosjek\_ctr =} \FunctionTok{round}\NormalTok{(}\FunctionTok{mean}\NormalTok{(click\_rate), }\DecValTok{4}\NormalTok{),}
      \AttributeTok{ukupni\_doseg =} \FunctionTok{sum}\NormalTok{(subscribers),}
      \AttributeTok{.groups =} \StringTok{"drop"}
\NormalTok{    )}
\NormalTok{\}}

\FunctionTok{sazetak\_kampanje}\NormalTok{(nl, }\StringTok{"weekly\_digest"}\NormalTok{)}
\end{Highlighting}
\end{Shaded}

\begin{verbatim}
# A tibble: 1 x 5
  tip               n prosjek_or prosjek_ctr ukupni_doseg
  <chr>         <int>      <dbl>       <dbl>        <dbl>
1 weekly_digest    12      0.289       0.029       217303
\end{verbatim}

\begin{Shaded}
\begin{Highlighting}[]
\FunctionTok{sazetak\_kampanje}\NormalTok{(nl, }\StringTok{"breaking\_news"}\NormalTok{)}
\end{Highlighting}
\end{Shaded}

\begin{verbatim}
# A tibble: 1 x 5
  tip               n prosjek_or prosjek_ctr ukupni_doseg
  <chr>         <int>      <dbl>       <dbl>        <dbl>
1 breaking_news     7      0.198      0.0382       120123
\end{verbatim}

Sad umjesto kopiranja pet blokova koda, pozivamo jednu funkciju s
različitim argumentom. Ako želite promijeniti izračun (dodati novu
metriku), mijenjate na jednom mjestu i promjena se automatski
primjenjuje svugdje.

\subsection{Funkcija koja vraća
graf}\label{funkcija-koja-vraux107a-graf}

Funkcije mogu vraćati bilo koji R objekt, uključujući ggplot grafove.

\begin{Shaded}
\begin{Highlighting}[]
\NormalTok{graf\_distribucije }\OtherTok{\textless{}{-}} \ControlFlowTok{function}\NormalTok{(data, varijabla, naslov) \{}
\NormalTok{  data }\SpecialCharTok{|\textgreater{}}
    \FunctionTok{ggplot}\NormalTok{(}\FunctionTok{aes}\NormalTok{(}\AttributeTok{x =}\NormalTok{ .data[[varijabla]])) }\SpecialCharTok{+}
    \FunctionTok{geom\_histogram}\NormalTok{(}\AttributeTok{fill =} \StringTok{"steelblue"}\NormalTok{, }\AttributeTok{color =} \StringTok{"white"}\NormalTok{, }\AttributeTok{bins =} \DecValTok{12}\NormalTok{) }\SpecialCharTok{+}
    \FunctionTok{labs}\NormalTok{(}\AttributeTok{title =}\NormalTok{ naslov, }\AttributeTok{x =}\NormalTok{ varijabla, }\AttributeTok{y =} \StringTok{"Broj kampanja"}\NormalTok{) }\SpecialCharTok{+}
    \FunctionTok{theme\_minimal}\NormalTok{()}
\NormalTok{\}}

\FunctionTok{graf\_distribucije}\NormalTok{(nl, }\StringTok{"open\_rate"}\NormalTok{, }\StringTok{"Distribucija open ratea"}\NormalTok{)}
\end{Highlighting}
\end{Shaded}

\pandocbounded{\includegraphics[keepaspectratio]{week-06_files/figure-pdf/funkcija-graf-1.pdf}}

Konstrukcija \texttt{.data{[}{[}varijabla{]}{]}} omogućuje
prosljeđivanje imena stupca kao teksta. Ovo je tehnički detalj tidyverse
programiranja koji je koristan kad pišete funkcije koje rade s
različitim stupcima.

\begin{Shaded}
\begin{Highlighting}[]
\FunctionTok{graf\_distribucije}\NormalTok{(nl, }\StringTok{"click\_rate"}\NormalTok{, }\StringTok{"Distribucija click ratea"}\NormalTok{)}
\end{Highlighting}
\end{Shaded}

\pandocbounded{\includegraphics[keepaspectratio]{week-06_files/figure-pdf/funkcija-graf-2-1.pdf}}

Ista funkcija, druga varijabla, novi graf. Ovo je suština DRY principa:
napišete logiku jednom i koristite je koliko god puta trebate.

\begin{tcolorbox}[enhanced jigsaw, coltitle=black, arc=.35mm, colframe=quarto-callout-tip-color-frame, left=2mm, bottomrule=.15mm, toprule=.15mm, leftrule=.75mm, bottomtitle=1mm, colback=white, title=\textcolor{quarto-callout-tip-color}{\faLightbulb}\hspace{0.5em}{Praktični savjet}, opacityback=0, colbacktitle=quarto-callout-tip-color!10!white, toptitle=1mm, titlerule=0mm, opacitybacktitle=0.6, breakable, rightrule=.15mm]

Pravilo tri: ako ste kopirali isti blok koda tri puta ili više, vrijeme
je da ga pretvorite u funkciju. Dva kopiranja su još prihvatljiva
(ponekad je brže kopirati nego pisati funkciju), ali tri signaliziraju
obrazac koji će se ponavljati i dalje. Funkcija vam štedi vrijeme
dugoročno i smanjuje rizik od grešaka pri kopiranju.

\end{tcolorbox}

\begin{center}\rule{0.5\linewidth}{0.5pt}\end{center}

\section{Uvjetne naredbe: if i else}\label{sec-uvjeti}

Uvjetne naredbe omogućuju R-u da donese odluku: ako je uvjet ispunjen,
napravi jedno, inače napravi drugo. Već smo koristili
\texttt{if\_else()} i \texttt{case\_when()} unutar \texttt{mutate()} za
rekodiranje varijabli. Sad učimo klasične \texttt{if}/\texttt{else}
naredbe koje rade izvan tibble konteksta.

\subsection{Osnovna sintaksa}\label{osnovna-sintaksa}

\begin{Shaded}
\begin{Highlighting}[]
\NormalTok{prosjek\_or }\OtherTok{\textless{}{-}} \FunctionTok{mean}\NormalTok{(nl}\SpecialCharTok{$}\NormalTok{open\_rate)}

\ControlFlowTok{if}\NormalTok{ (prosjek\_or }\SpecialCharTok{\textgreater{}} \FloatTok{0.25}\NormalTok{) \{}
  \FunctionTok{cat}\NormalTok{(}\StringTok{"Prosječni open rate je iznad 25\%, što je odličan rezultat.}\SpecialCharTok{\textbackslash{}n}\StringTok{"}\NormalTok{)}
\NormalTok{\} }\ControlFlowTok{else}\NormalTok{ \{}
  \FunctionTok{cat}\NormalTok{(}\StringTok{"Prosječni open rate je ispod 25\%, ima prostora za poboljšanje.}\SpecialCharTok{\textbackslash{}n}\StringTok{"}\NormalTok{)}
\NormalTok{\}}
\end{Highlighting}
\end{Shaded}

\begin{verbatim}
Prosječni open rate je iznad 25%, što je odličan rezultat.
\end{verbatim}

R evaluira uvjet u zagradi. Ako je TRUE, izvršava kod u prvom bloku. Ako
je FALSE, izvršava kod u else bloku. Funkcija \texttt{cat()} ispisuje
tekst u konzolu (slično \texttt{print()}, ali bez dodatnih oznaka).

\subsection{if, else if, else}\label{if-else-if-else}

Za više od dva ishoda, koristite \texttt{else\ if}.

\begin{Shaded}
\begin{Highlighting}[]
\NormalTok{ocijeni\_kampanju }\OtherTok{\textless{}{-}} \ControlFlowTok{function}\NormalTok{(open\_rate) \{}
  \ControlFlowTok{if}\NormalTok{ (open\_rate }\SpecialCharTok{\textgreater{}} \FloatTok{0.30}\NormalTok{) \{}
    \StringTok{"izvrsna"}
\NormalTok{  \} }\ControlFlowTok{else} \ControlFlowTok{if}\NormalTok{ (open\_rate }\SpecialCharTok{\textgreater{}} \FloatTok{0.20}\NormalTok{) \{}
    \StringTok{"dobra"}
\NormalTok{  \} }\ControlFlowTok{else} \ControlFlowTok{if}\NormalTok{ (open\_rate }\SpecialCharTok{\textgreater{}} \FloatTok{0.10}\NormalTok{) \{}
    \StringTok{"prosječna"}
\NormalTok{  \} }\ControlFlowTok{else}\NormalTok{ \{}
    \StringTok{"loša"}
\NormalTok{  \}}
\NormalTok{\}}

\FunctionTok{ocijeni\_kampanju}\NormalTok{(}\FloatTok{0.35}\NormalTok{)}
\end{Highlighting}
\end{Shaded}

\begin{verbatim}
[1] "izvrsna"
\end{verbatim}

\begin{Shaded}
\begin{Highlighting}[]
\FunctionTok{ocijeni\_kampanju}\NormalTok{(}\FloatTok{0.22}\NormalTok{)}
\end{Highlighting}
\end{Shaded}

\begin{verbatim}
[1] "dobra"
\end{verbatim}

\begin{Shaded}
\begin{Highlighting}[]
\FunctionTok{ocijeni\_kampanju}\NormalTok{(}\FloatTok{0.08}\NormalTok{)}
\end{Highlighting}
\end{Shaded}

\begin{verbatim}
[1] "loša"
\end{verbatim}

Uvjeti se provjeravaju redom, od vrha prema dnu. Čim je jedan uvjet
TRUE, pripadajuća vrijednost se vraća i R ne provjerava preostale
uvjete. Zato uvjete postavljamo od najstrožeg prema najblažem.

\subsection{Razlika između if/else i
if\_else()/case\_when()}\label{razlika-izmeux111u-ifelse-i-if_elsecase_when}

Ovo je česta točka zbunjenosti. Postoje dva različita sustava uvjetnog
izvršavanja u R-u i svaki ima svoje mjesto.

Klasični \texttt{if}/\texttt{else} radi s jednom vrijednošću. Koristi se
u funkcijama, skriptama i kontroli toka programa. Nije vektoriziran, što
znači da ne može obrađivati cijeli stupac odjednom.

\texttt{if\_else()} i \texttt{case\_when()} su vektorizirane funkcije.
Rade s cijelim vektorom (stupcem) odjednom i koriste se unutar
\texttt{mutate()} za rekodiranje varijabli u tibbleu.

\begin{Shaded}
\begin{Highlighting}[]
\CommentTok{\# if\_else() unutar mutate: radi na cijelom stupcu}
\NormalTok{nl }\SpecialCharTok{|\textgreater{}}
  \FunctionTok{mutate}\NormalTok{(}
    \AttributeTok{ocjena =} \FunctionTok{if\_else}\NormalTok{(open\_rate }\SpecialCharTok{\textgreater{}} \FloatTok{0.25}\NormalTok{, }\StringTok{"iznad prosjeka"}\NormalTok{, }\StringTok{"ispod prosjeka"}\NormalTok{)}
\NormalTok{  ) }\SpecialCharTok{|\textgreater{}}
  \FunctionTok{count}\NormalTok{(ocjena)}
\end{Highlighting}
\end{Shaded}

\begin{verbatim}
# A tibble: 2 x 2
  ocjena             n
  <chr>          <int>
1 ispod prosjeka    25
2 iznad prosjeka    25
\end{verbatim}

\begin{Shaded}
\begin{Highlighting}[]
\CommentTok{\# Klasični if/else: radi s jednom vrijednošću}
\CommentTok{\# (koristili smo ga u funkciji ocijeni\_kampanju)}
\end{Highlighting}
\end{Shaded}

Pravilo je jednostavno. Unutar \texttt{mutate()} koristite
\texttt{if\_else()} ili \texttt{case\_when()}. Izvan \texttt{mutate()},
u funkcijama i skriptama, koristite klasični \texttt{if}/\texttt{else}.

\subsection{Uvjeti u funkcijama: validacija
ulaza}\label{uvjeti-u-funkcijama-validacija-ulaza}

Praktična primjena \texttt{if}/\texttt{else} u funkcijama je provjera
jesu li ulazni podaci ispravni.

\begin{Shaded}
\begin{Highlighting}[]
\NormalTok{izracunaj\_ctr }\OtherTok{\textless{}{-}} \ControlFlowTok{function}\NormalTok{(clicks, impressions) \{}
  \ControlFlowTok{if}\NormalTok{ (impressions }\SpecialCharTok{\textless{}=} \DecValTok{0}\NormalTok{) \{}
    \FunctionTok{warning}\NormalTok{(}\StringTok{"Broj impresija mora biti pozitivan. Vraćam NA."}\NormalTok{)}
    \FunctionTok{return}\NormalTok{(}\ConstantTok{NA\_real\_}\NormalTok{)}
\NormalTok{  \}}
  
  \ControlFlowTok{if}\NormalTok{ (clicks }\SpecialCharTok{\textless{}} \DecValTok{0}\NormalTok{) \{}
    \FunctionTok{warning}\NormalTok{(}\StringTok{"Broj klikova ne može biti negativan. Vraćam NA."}\NormalTok{)}
    \FunctionTok{return}\NormalTok{(}\ConstantTok{NA\_real\_}\NormalTok{)}
\NormalTok{  \}}
  
  \FunctionTok{round}\NormalTok{(clicks }\SpecialCharTok{/}\NormalTok{ impressions, }\DecValTok{4}\NormalTok{)}
\NormalTok{\}}

\FunctionTok{izracunaj\_ctr}\NormalTok{(}\DecValTok{150}\NormalTok{, }\DecValTok{5000}\NormalTok{)}
\end{Highlighting}
\end{Shaded}

\begin{verbatim}
[1] 0.03
\end{verbatim}

\begin{Shaded}
\begin{Highlighting}[]
\FunctionTok{izracunaj\_ctr}\NormalTok{(}\DecValTok{150}\NormalTok{, }\DecValTok{0}\NormalTok{)}
\end{Highlighting}
\end{Shaded}

\begin{verbatim}
[1] NA
\end{verbatim}

\begin{Shaded}
\begin{Highlighting}[]
\FunctionTok{izracunaj\_ctr}\NormalTok{(}\SpecialCharTok{{-}}\DecValTok{10}\NormalTok{, }\DecValTok{5000}\NormalTok{)}
\end{Highlighting}
\end{Shaded}

\begin{verbatim}
[1] NA
\end{verbatim}

Funkcija \texttt{warning()} ispisuje upozorenje ali ne zaustavlja
izvršavanje. Funkcija \texttt{return()} eksplicitno vraća vrijednost i
izlazi iz funkcije. Bez \texttt{return()}, funkcija bi nastavila
izvršavanje i pokušala podijeliti s nulom.

Validacija ulaza je ono što razdvaja robusne funkcije od krhkih. Kad
pišete funkciju za sebe, možda znate da nikad nećete unijeti negativan
broj. Ali kad tu funkciju koristi netko drugi (ili vi za šest mjeseci,
kad ste zaboravili detalje), validacija sprečava tihe greške.

\begin{center}\rule{0.5\linewidth}{0.5pt}\end{center}

\section{For petlje: ponavljanje operacija}\label{sec-for}

Petlja je naredba koja ponavlja blok koda za svaki element u skupu.
\texttt{for} petlja u R-u ima jednostavnu sintaksu.

\subsection{Osnovna for petlja}\label{osnovna-for-petlja}

\begin{Shaded}
\begin{Highlighting}[]
\NormalTok{tipovi }\OtherTok{\textless{}{-}} \FunctionTok{unique}\NormalTok{(nl}\SpecialCharTok{$}\NormalTok{campaign\_type)}

\ControlFlowTok{for}\NormalTok{ (tip }\ControlFlowTok{in}\NormalTok{ tipovi) \{}
\NormalTok{  n }\OtherTok{\textless{}{-}}\NormalTok{ nl }\SpecialCharTok{|\textgreater{}} \FunctionTok{filter}\NormalTok{(campaign\_type }\SpecialCharTok{==}\NormalTok{ tip) }\SpecialCharTok{|\textgreater{}} \FunctionTok{nrow}\NormalTok{()}
  \FunctionTok{cat}\NormalTok{(tip, }\StringTok{":"}\NormalTok{, n, }\StringTok{"kampanja}\SpecialCharTok{\textbackslash{}n}\StringTok{"}\NormalTok{)}
\NormalTok{\}}
\end{Highlighting}
\end{Shaded}

\begin{verbatim}
special_report : 17 kampanja
weekly_digest : 12 kampanja
breaking_news : 7 kampanja
sponsored : 8 kampanja
event_promo : 6 kampanja
\end{verbatim}

R prolazi kroz svaki element vektora \texttt{tipovi}, dodjeljuje ga
varijabli \texttt{tip}, i izvršava kod u tijelu petlje. Kad se tijelo
izvrši za zadnji element, petlja završava.

\subsection{For petlja za generiranje
rezultata}\label{for-petlja-za-generiranje-rezultata}

Čest obrazac je korištenje petlje za prikupljanje rezultata u listu ili
tibble.

\begin{Shaded}
\begin{Highlighting}[]
\CommentTok{\# Inicijalizirajte praznu listu za rezultate}
\NormalTok{rezultati }\OtherTok{\textless{}{-}} \FunctionTok{list}\NormalTok{()}

\ControlFlowTok{for}\NormalTok{ (tip }\ControlFlowTok{in}\NormalTok{ tipovi) \{}
\NormalTok{  saz }\OtherTok{\textless{}{-}}\NormalTok{ nl }\SpecialCharTok{|\textgreater{}}
    \FunctionTok{filter}\NormalTok{(campaign\_type }\SpecialCharTok{==}\NormalTok{ tip) }\SpecialCharTok{|\textgreater{}}
    \FunctionTok{summarise}\NormalTok{(}
      \AttributeTok{tip =}\NormalTok{ tip,}
      \AttributeTok{n =} \FunctionTok{n}\NormalTok{(),}
      \AttributeTok{prosjek\_or =} \FunctionTok{round}\NormalTok{(}\FunctionTok{mean}\NormalTok{(open\_rate), }\DecValTok{3}\NormalTok{),}
      \AttributeTok{prosjek\_ctr =} \FunctionTok{round}\NormalTok{(}\FunctionTok{mean}\NormalTok{(click\_rate), }\DecValTok{4}\NormalTok{)}
\NormalTok{    )}
  
\NormalTok{  rezultati[[tip]] }\OtherTok{\textless{}{-}}\NormalTok{ saz}
\NormalTok{\}}

\CommentTok{\# Spojite sve rezultate u jedan tibble}
\FunctionTok{bind\_rows}\NormalTok{(rezultati)}
\end{Highlighting}
\end{Shaded}

\begin{verbatim}
# A tibble: 5 x 4
  tip                n prosjek_or prosjek_ctr
  <chr>          <int>      <dbl>       <dbl>
1 special_report    17      0.259      0.0519
2 weekly_digest     12      0.289      0.029 
3 breaking_news      7      0.198      0.0382
4 sponsored          8      0.248      0.0223
5 event_promo        6      0.25       0.0294
\end{verbatim}

Kreiramo praznu listu \texttt{rezultati}, u svakoj iteraciji računamo
sažetak i pohranjujemo ga u listu pod imenom tipa kampanje, a na kraju
sve spajamo u jedan tibble s \texttt{bind\_rows()}.

\subsection{For petlja za generiranje
grafova}\label{for-petlja-za-generiranje-grafova}

\begin{Shaded}
\begin{Highlighting}[]
\CommentTok{\# Generiranje grafa za svaki tip kampanje}
\ControlFlowTok{for}\NormalTok{ (tip }\ControlFlowTok{in} \FunctionTok{c}\NormalTok{(}\StringTok{"weekly\_digest"}\NormalTok{, }\StringTok{"breaking\_news"}\NormalTok{)) \{}
\NormalTok{  p }\OtherTok{\textless{}{-}}\NormalTok{ nl }\SpecialCharTok{|\textgreater{}}
    \FunctionTok{filter}\NormalTok{(campaign\_type }\SpecialCharTok{==}\NormalTok{ tip) }\SpecialCharTok{|\textgreater{}}
    \FunctionTok{ggplot}\NormalTok{(}\FunctionTok{aes}\NormalTok{(}\AttributeTok{x =}\NormalTok{ open\_rate)) }\SpecialCharTok{+}
    \FunctionTok{geom\_histogram}\NormalTok{(}\AttributeTok{fill =} \StringTok{"steelblue"}\NormalTok{, }\AttributeTok{color =} \StringTok{"white"}\NormalTok{, }\AttributeTok{bins =} \DecValTok{8}\NormalTok{) }\SpecialCharTok{+}
    \FunctionTok{labs}\NormalTok{(}
      \AttributeTok{title =} \FunctionTok{paste}\NormalTok{(}\StringTok{"Open rate:"}\NormalTok{, tip),}
      \AttributeTok{x =} \StringTok{"Open rate"}\NormalTok{,}
      \AttributeTok{y =} \StringTok{"Broj"}
\NormalTok{    ) }\SpecialCharTok{+}
    \FunctionTok{theme\_minimal}\NormalTok{()}
  
  \FunctionTok{print}\NormalTok{(p)}
\NormalTok{\}}
\end{Highlighting}
\end{Shaded}

\pandocbounded{\includegraphics[keepaspectratio]{week-06_files/figure-pdf/for-grafovi-1.pdf}}

\pandocbounded{\includegraphics[keepaspectratio]{week-06_files/figure-pdf/for-grafovi-2.pdf}}

Unutar for petlje, ggplot grafove morate eksplicitno ispisati s
\texttt{print()}. Izvan petlje, R automatski ispisuje zadnji objekt, ali
unutar petlje to ne radi. Ovo je čest izvor frustracije za početnike.

\begin{tcolorbox}[enhanced jigsaw, coltitle=black, arc=.35mm, colframe=quarto-callout-important-color-frame, left=2mm, bottomrule=.15mm, toprule=.15mm, leftrule=.75mm, bottomtitle=1mm, colback=white, title=\textcolor{quarto-callout-important-color}{\faExclamation}\hspace{0.5em}{Važna napomena}, opacityback=0, colbacktitle=quarto-callout-important-color!10!white, toptitle=1mm, titlerule=0mm, opacitybacktitle=0.6, breakable, rightrule=.15mm]

For petlje u R-u nisu pogrešne ni zastarjele, ali za većinu zadataka u
tidyverse ekosustavu postoje elegantnije alternative.
\texttt{group\_by()\ \textbar{}\textgreater{}\ summarise()} zamjenjuje
petlje za grupirane sažetke. \texttt{across()} zamjenjuje petlje za
primjenu iste operacije na više stupaca. \texttt{map()} iz paketa purrr
zamjenjuje petlje za primjenu funkcije na svaki element liste ili
vektora. Petlje koristite kad alternative ne postoje ili kad je petlja
jasnija (što se ponekad događa).

\end{tcolorbox}

\begin{center}\rule{0.5\linewidth}{0.5pt}\end{center}

\section{map(): moderna alternativa petljama}\label{sec-map}

Paket purrr (dio tidyverse) pruža obitelj \texttt{map()} funkcija koje
primjenjuju funkciju na svaki element vektora ili liste. Rezultat ovisi
o varijanti map-a koju koristite.

\subsection{Osnovni map()}\label{osnovni-map}

\begin{Shaded}
\begin{Highlighting}[]
\NormalTok{tipovi }\OtherTok{\textless{}{-}} \FunctionTok{unique}\NormalTok{(nl}\SpecialCharTok{$}\NormalTok{campaign\_type)}

\CommentTok{\# map() vraća listu}
\NormalTok{rezultati }\OtherTok{\textless{}{-}} \FunctionTok{map}\NormalTok{(tipovi, }\ControlFlowTok{function}\NormalTok{(tip) \{}
\NormalTok{  nl }\SpecialCharTok{|\textgreater{}}
    \FunctionTok{filter}\NormalTok{(campaign\_type }\SpecialCharTok{==}\NormalTok{ tip) }\SpecialCharTok{|\textgreater{}}
    \FunctionTok{summarise}\NormalTok{(}
      \AttributeTok{tip =}\NormalTok{ tip,}
      \AttributeTok{n =} \FunctionTok{n}\NormalTok{(),}
      \AttributeTok{prosjek\_or =} \FunctionTok{round}\NormalTok{(}\FunctionTok{mean}\NormalTok{(open\_rate), }\DecValTok{3}\NormalTok{)}
\NormalTok{    )}
\NormalTok{\})}

\FunctionTok{bind\_rows}\NormalTok{(rezultati)}
\end{Highlighting}
\end{Shaded}

\begin{verbatim}
# A tibble: 5 x 3
  tip                n prosjek_or
  <chr>          <int>      <dbl>
1 special_report    17      0.259
2 weekly_digest     12      0.289
3 breaking_news      7      0.198
4 sponsored          8      0.248
5 event_promo        6      0.25 
\end{verbatim}

\texttt{map()} prima vektor (ili listu) i funkciju, primjenjuje funkciju
na svaki element i vraća listu rezultata. Ovo je funkcionalni ekvivalent
for petlje ali u kompaktnijem obliku.

\subsection{Skraćena lambda
sintaksa}\label{skraux107ena-lambda-sintaksa}

Umjesto \texttt{function(tip)\ \{\ ...\ \}} možete koristiti skraćenu
lambda sintaksu s tildom.

\begin{Shaded}
\begin{Highlighting}[]
\CommentTok{\# Skraćena lambda sintaksa: \textbackslash{}(x) umjesto function(x)}
\NormalTok{rezultati }\OtherTok{\textless{}{-}} \FunctionTok{map}\NormalTok{(tipovi, \textbackslash{}(tip) \{}
\NormalTok{  nl }\SpecialCharTok{|\textgreater{}}
    \FunctionTok{filter}\NormalTok{(campaign\_type }\SpecialCharTok{==}\NormalTok{ tip) }\SpecialCharTok{|\textgreater{}}
    \FunctionTok{summarise}\NormalTok{(}\AttributeTok{tip =}\NormalTok{ tip, }\AttributeTok{n =} \FunctionTok{n}\NormalTok{(), }\AttributeTok{prosjek\_or =} \FunctionTok{round}\NormalTok{(}\FunctionTok{mean}\NormalTok{(open\_rate), }\DecValTok{3}\NormalTok{))}
\NormalTok{\})}

\FunctionTok{bind\_rows}\NormalTok{(rezultati)}
\end{Highlighting}
\end{Shaded}

\begin{verbatim}
# A tibble: 5 x 3
  tip                n prosjek_or
  <chr>          <int>      <dbl>
1 special_report    17      0.259
2 weekly_digest     12      0.289
3 breaking_news      7      0.198
4 sponsored          8      0.248
5 event_promo        6      0.25 
\end{verbatim}

Notacija \texttt{\textbackslash{}(tip)} je R-ova nova (od verzije 4.1)
skraćenica za \texttt{function(tip)}. Obje verzije rade identično, ali
\texttt{\textbackslash{}(x)} je kraća za pisanje.

\subsection{Varijante map-a}\label{varijante-map-a}

\texttt{map()} uvijek vraća listu. Kad znate kakav tip rezultata
očekujete, koristite specifičniju varijantu.

\begin{Shaded}
\begin{Highlighting}[]
\CommentTok{\# map\_dbl() vraća numerički vektor}
\NormalTok{prosjeci }\OtherTok{\textless{}{-}} \FunctionTok{map\_dbl}\NormalTok{(tipovi, \textbackslash{}(tip) \{}
\NormalTok{  nl }\SpecialCharTok{|\textgreater{}}
    \FunctionTok{filter}\NormalTok{(campaign\_type }\SpecialCharTok{==}\NormalTok{ tip) }\SpecialCharTok{|\textgreater{}}
    \FunctionTok{pull}\NormalTok{(open\_rate) }\SpecialCharTok{|\textgreater{}}
    \FunctionTok{mean}\NormalTok{()}
\NormalTok{\})}

\FunctionTok{tibble}\NormalTok{(}\AttributeTok{tip =}\NormalTok{ tipovi, }\AttributeTok{prosjek\_or =} \FunctionTok{round}\NormalTok{(prosjeci, }\DecValTok{3}\NormalTok{))}
\end{Highlighting}
\end{Shaded}

\begin{verbatim}
# A tibble: 5 x 2
  tip            prosjek_or
  <chr>               <dbl>
1 special_report      0.259
2 weekly_digest       0.289
3 breaking_news       0.198
4 sponsored           0.248
5 event_promo         0.25 
\end{verbatim}

\texttt{map\_dbl()} vraća numerički (double) vektor umjesto liste.
\texttt{map\_chr()} vraća tekstualni vektor. \texttt{map\_lgl()} vraća
logički. \texttt{map\_df()} vraća tibble (spaja sve rezultate).
Koristite specifičnu varijantu kad god možete jer je rezultat
predvidljiviji i jednostavniji za daljnji rad.

\subsection{map() unutar tibble radnog
toka}\label{map-unutar-tibble-radnog-toka}

Najelegantnija primjena \texttt{map()} je unutar tibble radnog toka,
kombinirano s \texttt{nest()} i \texttt{unnest()}.

\begin{Shaded}
\begin{Highlighting}[]
\NormalTok{nl }\SpecialCharTok{|\textgreater{}}
  \FunctionTok{group\_by}\NormalTok{(campaign\_type) }\SpecialCharTok{|\textgreater{}}
  \FunctionTok{nest}\NormalTok{() }\SpecialCharTok{|\textgreater{}}
  \FunctionTok{mutate}\NormalTok{(}
    \AttributeTok{n =} \FunctionTok{map\_dbl}\NormalTok{(data, nrow),}
    \AttributeTok{prosjek\_or =} \FunctionTok{map\_dbl}\NormalTok{(data, \textbackslash{}(df) }\FunctionTok{mean}\NormalTok{(df}\SpecialCharTok{$}\NormalTok{open\_rate)),}
    \AttributeTok{prosjek\_ctr =} \FunctionTok{map\_dbl}\NormalTok{(data, \textbackslash{}(df) }\FunctionTok{mean}\NormalTok{(df}\SpecialCharTok{$}\NormalTok{click\_rate))}
\NormalTok{  ) }\SpecialCharTok{|\textgreater{}}
  \FunctionTok{select}\NormalTok{(}\SpecialCharTok{{-}}\NormalTok{data) }\SpecialCharTok{|\textgreater{}}
  \FunctionTok{mutate}\NormalTok{(}\FunctionTok{across}\NormalTok{(}\FunctionTok{c}\NormalTok{(prosjek\_or, prosjek\_ctr), \textbackslash{}(x) }\FunctionTok{round}\NormalTok{(x, }\DecValTok{4}\NormalTok{)))}
\end{Highlighting}
\end{Shaded}

\begin{verbatim}
# A tibble: 5 x 4
# Groups:   campaign_type [5]
  campaign_type      n prosjek_or prosjek_ctr
  <chr>          <dbl>      <dbl>       <dbl>
1 special_report    17      0.259      0.0519
2 weekly_digest     12      0.289      0.029 
3 breaking_news      7      0.198      0.0382
4 sponsored          8      0.248      0.0223
5 event_promo        6      0.250      0.0294
\end{verbatim}

Funkcija \texttt{nest()} pakira podatke svake grupe u zasebni tibble
unutar liste-stupca \texttt{data}. Zatim \texttt{map\_dbl()} primjenjuje
funkciju na svaki od tih ugniježđenih tibbleova. Rezultat je jedan redak
po grupi s izračunatim metrikama.

Ovo je napredni obrazac koji ćete cijeniti kad budete radili složenije
analize (na primjer, fitanje zasebnog regresijskog modela za svaku grupu
u tjednu 14).

\begin{center}\rule{0.5\linewidth}{0.5pt}\end{center}

\section{DRY princip i organizacija skripte}\label{sec-dry}

DRY (Don't Repeat Yourself) je programerski princip koji kaže da svaka
informacija u kodu treba postojati na jednom mjestu. Kad se ponavljate,
stvarate više točaka koje trebate ažurirati kad nešto promijenite, a to
je recept za greške.

\subsection{Primjer: parametri na jednom
mjestu}\label{primjer-parametri-na-jednom-mjestu}

\begin{Shaded}
\begin{Highlighting}[]
\CommentTok{\# PARAMETRI (mijenjajte ovdje, promjena se propagira svugdje)}
\NormalTok{min\_kampanja\_za\_analizu }\OtherTok{\textless{}{-}} \DecValTok{5}
\NormalTok{decimale }\OtherTok{\textless{}{-}} \DecValTok{3}
\NormalTok{boja\_grafova }\OtherTok{\textless{}{-}} \StringTok{"steelblue"}
\NormalTok{kategorije\_interesa }\OtherTok{\textless{}{-}} \FunctionTok{c}\NormalTok{(}\StringTok{"weekly\_digest"}\NormalTok{, }\StringTok{"breaking\_news"}\NormalTok{, }\StringTok{"special\_report"}\NormalTok{)}

\CommentTok{\# ANALIZA (koristi parametre odozgo)}
\NormalTok{nl\_filtered }\OtherTok{\textless{}{-}}\NormalTok{ nl }\SpecialCharTok{|\textgreater{}}
  \FunctionTok{filter}\NormalTok{(campaign\_type }\SpecialCharTok{\%in\%}\NormalTok{ kategorije\_interesa)}

\NormalTok{nl\_filtered }\SpecialCharTok{|\textgreater{}}
  \FunctionTok{group\_by}\NormalTok{(campaign\_type) }\SpecialCharTok{|\textgreater{}}
  \FunctionTok{summarise}\NormalTok{(}
    \AttributeTok{n =} \FunctionTok{n}\NormalTok{(),}
    \AttributeTok{prosjek\_or =} \FunctionTok{round}\NormalTok{(}\FunctionTok{mean}\NormalTok{(open\_rate), decimale),}
    \AttributeTok{prosjek\_ctr =} \FunctionTok{round}\NormalTok{(}\FunctionTok{mean}\NormalTok{(click\_rate), decimale),}
    \AttributeTok{.groups =} \StringTok{"drop"}
\NormalTok{  ) }\SpecialCharTok{|\textgreater{}}
  \FunctionTok{filter}\NormalTok{(n }\SpecialCharTok{\textgreater{}=}\NormalTok{ min\_kampanja\_za\_analizu)}
\end{Highlighting}
\end{Shaded}

\begin{verbatim}
# A tibble: 3 x 4
  campaign_type      n prosjek_or prosjek_ctr
  <chr>          <int>      <dbl>       <dbl>
1 breaking_news      7      0.198       0.038
2 special_report    17      0.259       0.052
3 weekly_digest     12      0.289       0.029
\end{verbatim}

Svi ključni parametri su definirani na jednom mjestu na vrhu. Kad
klijent kaže ``pokaži mi analizu samo za weekly digest i special
report'', mijenjate jednu liniju i cijela analiza se ažurira. Ovo je
fundamentalno drugačije od traženja i zamjenjivanja vrijednosti
razbacanih po cijelom kodu.

\subsection{Struktura analitičke
skripte}\label{struktura-analitiux10dke-skripte}

Dobro organizirana skripta ima jasne sekcije. Svaka sekcija radi jednu
stvar i jasno je označena.

\begin{Shaded}
\begin{Highlighting}[]
\CommentTok{\# ============================================}
\CommentTok{\# Analiza newsletter kampanja}
\CommentTok{\# Autor: Ime Prezime}
\CommentTok{\# Datum: 2025{-}03{-}29}
\CommentTok{\# Opis: Sažetak performansi email kampanja}
\CommentTok{\# ============================================}

\CommentTok{\# 1. PAKETI {-}{-}{-}{-}}
\FunctionTok{library}\NormalTok{(tidyverse)}

\CommentTok{\# 2. PARAMETRI {-}{-}{-}{-}}
\NormalTok{input\_file }\OtherTok{\textless{}{-}} \StringTok{"../resources/datasets/newsletter\_campaigns.csv"}
\NormalTok{output\_dir }\OtherTok{\textless{}{-}} \StringTok{"../outputs/"}
\NormalTok{min\_n }\OtherTok{\textless{}{-}} \DecValTok{5}

\CommentTok{\# 3. UČITAVANJE {-}{-}{-}{-}}
\NormalTok{raw }\OtherTok{\textless{}{-}} \FunctionTok{read\_csv}\NormalTok{(input\_file)}

\CommentTok{\# 4. ČIŠĆENJE {-}{-}{-}{-}}
\NormalTok{clean }\OtherTok{\textless{}{-}}\NormalTok{ raw }\SpecialCharTok{|\textgreater{}}
  \FunctionTok{filter}\NormalTok{(}\SpecialCharTok{!}\FunctionTok{is.na}\NormalTok{(open\_rate)) }\SpecialCharTok{|\textgreater{}}
  \FunctionTok{mutate}\NormalTok{(}
    \AttributeTok{campaign\_type =} \FunctionTok{factor}\NormalTok{(campaign\_type),}
    \AttributeTok{ocjena =} \FunctionTok{case\_when}\NormalTok{(}
\NormalTok{      open\_rate }\SpecialCharTok{\textgreater{}} \FloatTok{0.30} \SpecialCharTok{\textasciitilde{}} \StringTok{"izvrsna"}\NormalTok{,}
\NormalTok{      open\_rate }\SpecialCharTok{\textgreater{}} \FloatTok{0.20} \SpecialCharTok{\textasciitilde{}} \StringTok{"dobra"}\NormalTok{,}
\NormalTok{      open\_rate }\SpecialCharTok{\textgreater{}} \FloatTok{0.10} \SpecialCharTok{\textasciitilde{}} \StringTok{"prosjecna"}\NormalTok{,}
      \AttributeTok{.default =} \StringTok{"losa"}
\NormalTok{    )}
\NormalTok{  )}

\CommentTok{\# 5. ANALIZA {-}{-}{-}{-}}
\NormalTok{sazetak }\OtherTok{\textless{}{-}}\NormalTok{ clean }\SpecialCharTok{|\textgreater{}}
  \FunctionTok{group\_by}\NormalTok{(campaign\_type) }\SpecialCharTok{|\textgreater{}}
  \FunctionTok{summarise}\NormalTok{(}
    \AttributeTok{n =} \FunctionTok{n}\NormalTok{(),}
    \AttributeTok{M\_or =} \FunctionTok{mean}\NormalTok{(open\_rate),}
    \AttributeTok{SD\_or =} \FunctionTok{sd}\NormalTok{(open\_rate),}
    \AttributeTok{M\_ctr =} \FunctionTok{mean}\NormalTok{(click\_rate),}
    \AttributeTok{.groups =} \StringTok{"drop"}
\NormalTok{  )}

\CommentTok{\# 6. VIZUALIZACIJA {-}{-}{-}{-}}
\NormalTok{graf }\OtherTok{\textless{}{-}} \FunctionTok{ggplot}\NormalTok{(clean, }\FunctionTok{aes}\NormalTok{(}\AttributeTok{x =}\NormalTok{ campaign\_type, }\AttributeTok{y =}\NormalTok{ open\_rate)) }\SpecialCharTok{+}
  \FunctionTok{geom\_boxplot}\NormalTok{(}\AttributeTok{fill =} \StringTok{"steelblue"}\NormalTok{, }\AttributeTok{alpha =} \FloatTok{0.6}\NormalTok{) }\SpecialCharTok{+}
  \FunctionTok{theme\_minimal}\NormalTok{() }\SpecialCharTok{+}
  \FunctionTok{labs}\NormalTok{(}\AttributeTok{title =} \StringTok{"Open rate po tipu kampanje"}\NormalTok{)}

\CommentTok{\# 7. IZVOZ {-}{-}{-}{-}}
\FunctionTok{write\_csv}\NormalTok{(sazetak, }\FunctionTok{paste0}\NormalTok{(output\_dir, }\StringTok{"sazetak\_kampanja.csv"}\NormalTok{))}
\FunctionTok{ggsave}\NormalTok{(}\FunctionTok{paste0}\NormalTok{(output\_dir, }\StringTok{"boxplot\_open\_rate.png"}\NormalTok{), graf, }\AttributeTok{width =} \DecValTok{8}\NormalTok{, }\AttributeTok{height =} \DecValTok{5}\NormalTok{)}
\end{Highlighting}
\end{Shaded}

Komentari s četiri crtice (\texttt{\#\ 1.\ PAKETI\ -\/-\/-\/-}) stvaraju
navigacijske oznake u Positronu (ili RStudiju) koje omogućuju brzo
skakanje između sekcija. Ovo je konvencija, ne sintaktičko pravilo, ali
je široko prihvaćena u R zajednici.

\begin{quote}
Vaša skripta je vaš laboratorijski dnevnik. Svaki korak je dokumentiran,
svaka odluka komentirana, svaki rezultat ponovljiv. Netko (uključujući
vas za šest mjeseci) mora moći pokrenuti skriptu od početka do kraja i
dobiti identične rezultate.
\end{quote}

\subsection{Pomoćne funkcije na vrhu
skripte}\label{pomoux107ne-funkcije-na-vrhu-skripte}

Kad imate funkcije koje koristite na više mjesta u analizi, definirajte
ih odmah nakon učitavanja paketa. Ovo ih čini vidljivima kroz cijelu
skriptu.

\begin{Shaded}
\begin{Highlighting}[]
\CommentTok{\# Pomoćne funkcije za newsletter analizu}
\NormalTok{sazetak\_metrike }\OtherTok{\textless{}{-}} \ControlFlowTok{function}\NormalTok{(data, metrika, }\AttributeTok{decimale =} \DecValTok{3}\NormalTok{) \{}
\NormalTok{  data }\SpecialCharTok{|\textgreater{}}
    \FunctionTok{summarise}\NormalTok{(}
      \AttributeTok{M =} \FunctionTok{round}\NormalTok{(}\FunctionTok{mean}\NormalTok{(.data[[metrika]], }\AttributeTok{na.rm =} \ConstantTok{TRUE}\NormalTok{), decimale),}
      \AttributeTok{Med =} \FunctionTok{round}\NormalTok{(}\FunctionTok{median}\NormalTok{(.data[[metrika]], }\AttributeTok{na.rm =} \ConstantTok{TRUE}\NormalTok{), decimale),}
      \AttributeTok{SD =} \FunctionTok{round}\NormalTok{(}\FunctionTok{sd}\NormalTok{(.data[[metrika]], }\AttributeTok{na.rm =} \ConstantTok{TRUE}\NormalTok{), decimale),}
      \AttributeTok{Min =} \FunctionTok{round}\NormalTok{(}\FunctionTok{min}\NormalTok{(.data[[metrika]], }\AttributeTok{na.rm =} \ConstantTok{TRUE}\NormalTok{), decimale),}
      \AttributeTok{Max =} \FunctionTok{round}\NormalTok{(}\FunctionTok{max}\NormalTok{(.data[[metrika]], }\AttributeTok{na.rm =} \ConstantTok{TRUE}\NormalTok{), decimale)}
\NormalTok{    )}
\NormalTok{\}}

\NormalTok{ocjena\_kampanje }\OtherTok{\textless{}{-}} \ControlFlowTok{function}\NormalTok{(open\_rate) \{}
  \FunctionTok{case\_when}\NormalTok{(}
\NormalTok{    open\_rate }\SpecialCharTok{\textgreater{}} \FloatTok{0.30} \SpecialCharTok{\textasciitilde{}} \StringTok{"izvrsna"}\NormalTok{,}
\NormalTok{    open\_rate }\SpecialCharTok{\textgreater{}} \FloatTok{0.20} \SpecialCharTok{\textasciitilde{}} \StringTok{"dobra"}\NormalTok{,}
\NormalTok{    open\_rate }\SpecialCharTok{\textgreater{}} \FloatTok{0.10} \SpecialCharTok{\textasciitilde{}} \StringTok{"prosjecna"}\NormalTok{,}
    \AttributeTok{.default =} \StringTok{"losa"}
\NormalTok{  )}
\NormalTok{\}}

\CommentTok{\# Korištenje pomoćnih funkcija}
\NormalTok{nl }\SpecialCharTok{|\textgreater{}}
  \FunctionTok{group\_by}\NormalTok{(campaign\_type) }\SpecialCharTok{|\textgreater{}}
  \FunctionTok{sazetak\_metrike}\NormalTok{(}\StringTok{"open\_rate"}\NormalTok{)}
\end{Highlighting}
\end{Shaded}

\begin{verbatim}
# A tibble: 5 x 6
  campaign_type      M   Med    SD   Min   Max
  <chr>          <dbl> <dbl> <dbl> <dbl> <dbl>
1 breaking_news  0.198 0.195 0.055 0.115 0.281
2 event_promo    0.25  0.23  0.086 0.161 0.408
3 special_report 0.259 0.255 0.067 0.155 0.393
4 sponsored      0.248 0.246 0.043 0.191 0.309
5 weekly_digest  0.289 0.29  0.066 0.178 0.388
\end{verbatim}

\begin{Shaded}
\begin{Highlighting}[]
\NormalTok{nl }\SpecialCharTok{|\textgreater{}}
  \FunctionTok{mutate}\NormalTok{(}\AttributeTok{ocjena =} \FunctionTok{ocjena\_kampanje}\NormalTok{(open\_rate)) }\SpecialCharTok{|\textgreater{}}
  \FunctionTok{count}\NormalTok{(ocjena, }\AttributeTok{sort =} \ConstantTok{TRUE}\NormalTok{)}
\end{Highlighting}
\end{Shaded}

\begin{verbatim}
# A tibble: 3 x 2
  ocjena        n
  <chr>     <int>
1 dobra        25
2 prosjecna    14
3 izvrsna      11
\end{verbatim}

Definirajući \texttt{ocjena\_kampanje()} kao funkciju, logiku
rekodiranja pišete jednom. Ako se kriteriji promijene (recimo, prag za
``izvrsno'' padne na 0.28), mijenjate na jednom mjestu.

\begin{center}\rule{0.5\linewidth}{0.5pt}\end{center}

\section{Praktični primjer: automatizirana analiza po
kampanjama}\label{sec-prakticni-primjer}

Spojimo sve naučene koncepte u jednom praktičnom primjeru. Cilj je
napisati kod koji za svaki tip kampanje generira sažetak tablica i graf,
koristeći funkcije, map i DRY principe.

\begin{Shaded}
\begin{Highlighting}[]
\CommentTok{\# Funkcija za kompletnu analizu jednog tipa kampanje}
\NormalTok{analiziraj\_tip }\OtherTok{\textless{}{-}} \ControlFlowTok{function}\NormalTok{(data, tip) \{}
\NormalTok{  podaci }\OtherTok{\textless{}{-}}\NormalTok{ data }\SpecialCharTok{|\textgreater{}} \FunctionTok{filter}\NormalTok{(campaign\_type }\SpecialCharTok{==}\NormalTok{ tip)}
  
  \ControlFlowTok{if}\NormalTok{ (}\FunctionTok{nrow}\NormalTok{(podaci) }\SpecialCharTok{\textless{}} \DecValTok{3}\NormalTok{) \{}
    \FunctionTok{return}\NormalTok{(}\ConstantTok{NULL}\NormalTok{)  }\CommentTok{\# Preskoči tipove s premalo podataka}
\NormalTok{  \}}
  
\NormalTok{  saz }\OtherTok{\textless{}{-}}\NormalTok{ podaci }\SpecialCharTok{|\textgreater{}}
    \FunctionTok{summarise}\NormalTok{(}
      \AttributeTok{tip =}\NormalTok{ tip,}
      \AttributeTok{n =} \FunctionTok{n}\NormalTok{(),}
      \AttributeTok{or\_M =} \FunctionTok{round}\NormalTok{(}\FunctionTok{mean}\NormalTok{(open\_rate), }\DecValTok{3}\NormalTok{),}
      \AttributeTok{or\_SD =} \FunctionTok{round}\NormalTok{(}\FunctionTok{sd}\NormalTok{(open\_rate), }\DecValTok{3}\NormalTok{),}
      \AttributeTok{ctr\_M =} \FunctionTok{round}\NormalTok{(}\FunctionTok{mean}\NormalTok{(click\_rate), }\DecValTok{4}\NormalTok{),}
      \AttributeTok{prosj\_pretplatnika =} \FunctionTok{round}\NormalTok{(}\FunctionTok{mean}\NormalTok{(subscribers), }\DecValTok{0}\NormalTok{),}
      \AttributeTok{prosj\_rijeci =} \FunctionTok{round}\NormalTok{(}\FunctionTok{mean}\NormalTok{(word\_count), }\DecValTok{0}\NormalTok{)}
\NormalTok{    )}
  
\NormalTok{  saz}
\NormalTok{\}}

\CommentTok{\# Primjena na sve tipove}
\NormalTok{svi\_tipovi }\OtherTok{\textless{}{-}} \FunctionTok{unique}\NormalTok{(nl}\SpecialCharTok{$}\NormalTok{campaign\_type)}

\NormalTok{rezultati }\OtherTok{\textless{}{-}} \FunctionTok{map}\NormalTok{(svi\_tipovi, \textbackslash{}(tip) }\FunctionTok{analiziraj\_tip}\NormalTok{(nl, tip)) }\SpecialCharTok{|\textgreater{}}
  \FunctionTok{bind\_rows}\NormalTok{()}

\NormalTok{rezultati }\SpecialCharTok{|\textgreater{}}
  \FunctionTok{arrange}\NormalTok{(}\FunctionTok{desc}\NormalTok{(or\_M))}
\end{Highlighting}
\end{Shaded}

\begin{verbatim}
# A tibble: 5 x 7
  tip                n  or_M or_SD  ctr_M prosj_pretplatnika prosj_rijeci
  <chr>          <int> <dbl> <dbl>  <dbl>              <dbl>        <dbl>
1 weekly_digest     12 0.289 0.066 0.029               18109          380
2 special_report    17 0.259 0.067 0.0519              17581          486
3 event_promo        6 0.25  0.086 0.0294              12938          210
4 sponsored          8 0.248 0.043 0.0223              13609          204
5 breaking_news      7 0.198 0.055 0.0382              17160          140
\end{verbatim}

Ovaj pristup ima nekoliko prednosti. Logika analize je enkapsulirana u
jednu funkciju. Validacija (\texttt{if\ (nrow(podaci)\ \textless{}\ 3)})
osigurava da ne radimo besmislene izračune na premalo podataka.
\texttt{map()} elegantno primjenjuje funkciju na sve tipove. Rezultat je
čist tibble sortiran po open rateu.

\begin{Shaded}
\begin{Highlighting}[]
\CommentTok{\# Vizualizacija rezultata}
\NormalTok{rezultati }\SpecialCharTok{|\textgreater{}}
  \FunctionTok{mutate}\NormalTok{(}\AttributeTok{tip =} \FunctionTok{fct\_reorder}\NormalTok{(tip, or\_M)) }\SpecialCharTok{|\textgreater{}}
  \FunctionTok{ggplot}\NormalTok{(}\FunctionTok{aes}\NormalTok{(}\AttributeTok{x =}\NormalTok{ tip, }\AttributeTok{y =}\NormalTok{ or\_M)) }\SpecialCharTok{+}
  \FunctionTok{geom\_col}\NormalTok{(}\AttributeTok{fill =} \StringTok{"steelblue"}\NormalTok{) }\SpecialCharTok{+}
  \FunctionTok{geom\_errorbar}\NormalTok{(}
    \FunctionTok{aes}\NormalTok{(}\AttributeTok{ymin =}\NormalTok{ or\_M }\SpecialCharTok{{-}}\NormalTok{ or\_SD, }\AttributeTok{ymax =}\NormalTok{ or\_M }\SpecialCharTok{+}\NormalTok{ or\_SD),}
    \AttributeTok{width =} \FloatTok{0.2}
\NormalTok{  ) }\SpecialCharTok{+}
  \FunctionTok{coord\_flip}\NormalTok{() }\SpecialCharTok{+}
  \FunctionTok{labs}\NormalTok{(}
    \AttributeTok{title =} \StringTok{"Prosječni open rate po tipu newsletter kampanje"}\NormalTok{,}
    \AttributeTok{subtitle =} \StringTok{"Stupac = prosjek, error bar = +/{-} 1 SD"}\NormalTok{,}
    \AttributeTok{x =} \ConstantTok{NULL}\NormalTok{,}
    \AttributeTok{y =} \StringTok{"Open rate"}
\NormalTok{  ) }\SpecialCharTok{+}
  \FunctionTok{theme\_minimal}\NormalTok{()}
\end{Highlighting}
\end{Shaded}

\pandocbounded{\includegraphics[keepaspectratio]{week-06_files/figure-pdf/prakticni-primjer-graf-1.pdf}}

Error barovi (crte pogreške) prikazuju jednu standardnu devijaciju iznad
i ispod prosjeka, dajući vizualni uvid u varijabilnost unutar svake
kategorije. Breaking news kampanje imaju viši prosječni open rate ali i
veću varijabilnost, dok su sponsored kampanje konzistentno niže.

\begin{center}\rule{0.5\linewidth}{0.5pt}\end{center}

\begin{tcolorbox}[enhanced jigsaw, coltitle=black, arc=.35mm, colframe=quarto-callout-note-color-frame, left=2mm, bottomrule=.15mm, toprule=.15mm, leftrule=.75mm, bottomtitle=1mm, colback=white, title=\textcolor{quarto-callout-note-color}{\faInfo}\hspace{0.5em}{Podsjetnik}, opacityback=0, colbacktitle=quarto-callout-note-color!10!white, toptitle=1mm, titlerule=0mm, opacitybacktitle=0.6, breakable, rightrule=.15mm]

U prvom dijelu naučili smo pisati vlastite funkcije, koristiti uvjetne
naredbe, for petlje i map() iz paketa purrr. U ovom dijelu primjenjujemo
te vještine na realne radne tokove: rad s više datoteka, debugging,
Quarto izvještaji i kompletna automatizirana analiza.

\end{tcolorbox}

\section{Rad s više datoteka}\label{sec-vise-datoteka}

U praksi, podaci rijetko dolaze u jednoj datoteci. Možda imate zasebne
CSV datoteke za svaki mjesec, za svaku kampanju ili za svaki izvor
podataka. Umjesto ručnog učitavanja svake datoteke, možete
automatizirati taj proces koristeći koncepte koje smo upravo naučili.

\subsection{Pronalaženje datoteka}\label{pronalaux17eenje-datoteka}

Funkcija \texttt{list.files()} pronalazi datoteke u direktoriju po
zadanom uzorku.

\begin{Shaded}
\begin{Highlighting}[]
\CommentTok{\# Popis svih CSV datoteka u datasets direktoriju}
\NormalTok{csv\_datoteke }\OtherTok{\textless{}{-}} \FunctionTok{list.files}\NormalTok{(}
  \AttributeTok{path =} \StringTok{"../resources/datasets/"}\NormalTok{,}
  \AttributeTok{pattern =} \StringTok{"}\SpecialCharTok{\textbackslash{}\textbackslash{}}\StringTok{.csv$"}\NormalTok{,}
  \AttributeTok{full.names =} \ConstantTok{TRUE}
\NormalTok{)}

\NormalTok{csv\_datoteke}
\end{Highlighting}
\end{Shaded}

\begin{verbatim}
 [1] "../resources/datasets/ab_test_headlines.csv"   
 [2] "../resources/datasets/article_engagement.csv"  
 [3] "../resources/datasets/instagram_ab_test.csv"   
 [4] "../resources/datasets/media_habits_raw.csv"    
 [5] "../resources/datasets/media_population.csv"    
 [6] "../resources/datasets/media_trust.csv"         
 [7] "../resources/datasets/newsletter_campaign.csv" 
 [8] "../resources/datasets/newsletter_campaigns.csv"
 [9] "../resources/datasets/social_media_survey.csv" 
[10] "../resources/datasets/social_posts.csv"        
[11] "../resources/datasets/tiktok_usage.csv"        
\end{verbatim}

Argument \texttt{pattern\ =\ "\textbackslash{}\textbackslash{}.csv\$"}
koristi regularni izraz za pronalaženje datoteka koje završavaju s .csv.
\texttt{full.names\ =\ TRUE} vraća kompletne putanje (ne samo imena
datoteka), što je bitno jer ih trebamo za učitavanje.

\subsection{Učitavanje više datoteka
odjednom}\label{uux10ditavanje-viux161e-datoteka-odjednom}

Kombinirajmo \texttt{list.files()}, \texttt{map()} i
\texttt{bind\_rows()} za učitavanje i spajanje svih CSV datoteka u
jednom koraku.

\begin{Shaded}
\begin{Highlighting}[]
\CommentTok{\# Učitaj sve CSV datoteke i spoji ih}
\NormalTok{svi\_podaci }\OtherTok{\textless{}{-}}\NormalTok{ csv\_datoteke }\SpecialCharTok{|\textgreater{}}
  \FunctionTok{map}\NormalTok{(\textbackslash{}(f) }\FunctionTok{read\_csv}\NormalTok{(f, }\AttributeTok{show\_col\_types =} \ConstantTok{FALSE}\NormalTok{)) }\SpecialCharTok{|\textgreater{}}
  \FunctionTok{bind\_rows}\NormalTok{()}
\end{Highlighting}
\end{Shaded}

Ovo je moćan obrazac. \texttt{map()} primjenjuje \texttt{read\_csv()} na
svaku putanju, vraćajući listu tibbleova. \texttt{bind\_rows()} ih
vertikalno spaja u jedan veliki tibble. Ako datoteke imaju iste stupce,
rezultat je jednostavna konkatenacija. Ako se stupci razlikuju,
\texttt{bind\_rows()} popunjava nedostajuće s NA.

\subsection{Dodavanje informacije o
izvoru}\label{dodavanje-informacije-o-izvoru}

Često želite znati iz koje datoteke dolazi koji redak. Funkcija
\texttt{set\_names()} pomaže.

\begin{Shaded}
\begin{Highlighting}[]
\CommentTok{\# Učitaj sve datoteke i dodaj stupac s imenom datoteke}
\NormalTok{svi\_podaci }\OtherTok{\textless{}{-}}\NormalTok{ csv\_datoteke }\SpecialCharTok{|\textgreater{}}
  \FunctionTok{set\_names}\NormalTok{() }\SpecialCharTok{|\textgreater{}}
  \FunctionTok{map}\NormalTok{(\textbackslash{}(f) }\FunctionTok{read\_csv}\NormalTok{(f, }\AttributeTok{show\_col\_types =} \ConstantTok{FALSE}\NormalTok{)) }\SpecialCharTok{|\textgreater{}}
  \FunctionTok{bind\_rows}\NormalTok{(}\AttributeTok{.id =} \StringTok{"izvor"}\NormalTok{)}
\end{Highlighting}
\end{Shaded}

Argument \texttt{.id\ =\ "izvor"} u \texttt{bind\_rows()} kreira novi
stupac \texttt{izvor} koji sadrži ime elementa liste (u ovom slučaju
putanju datoteke). Ovo je korisno za praćenje porijekla podataka.

\begin{tcolorbox}[enhanced jigsaw, coltitle=black, arc=.35mm, colframe=quarto-callout-tip-color-frame, left=2mm, bottomrule=.15mm, toprule=.15mm, leftrule=.75mm, bottomtitle=1mm, colback=white, title=\textcolor{quarto-callout-tip-color}{\faLightbulb}\hspace{0.5em}{Praktični savjet}, opacityback=0, colbacktitle=quarto-callout-tip-color!10!white, toptitle=1mm, titlerule=0mm, opacitybacktitle=0.6, breakable, rightrule=.15mm]

Obrazac
\texttt{list.files()\ \textbar{}\textgreater{}\ map(read\_csv)\ \textbar{}\textgreater{}\ bind\_rows()}
je jedan od najkorisnijih obrazaca u cijelom R radnom toku. Naučite ga
napamet. Koristit ćete ga svaki put kad dobijete podatke razdijeljene u
više datoteka (mjesečni izvještaji, odvojene ankete, logovi po danima).

\end{tcolorbox}

\begin{center}\rule{0.5\linewidth}{0.5pt}\end{center}

\section{Debugging: pronalaženje i ispravljanje
grešaka}\label{sec-debugging}

Greške su neizbježan dio programiranja. Pitanje nije hoćete li naletjeti
na grešku, nego koliko ćete brzo identificirati i ispraviti problem. R
daje poruke o greškama koje su ponekad jasne, a ponekad kriptične. Evo
strategija za sustavno traženje problema.

\subsection{Čitanje poruka o
greškama}\label{ux10ditanje-poruka-o-greux161kama}

\begin{Shaded}
\begin{Highlighting}[]
\CommentTok{\# Tipična greška: objekt ne postoji}
\NormalTok{nl }\SpecialCharTok{|\textgreater{}}
  \FunctionTok{filter}\NormalTok{(kampanja\_tip }\SpecialCharTok{==} \StringTok{"weekly\_digest"}\NormalTok{)}
\CommentTok{\# Error: object \textquotesingle{}kampanja\_tip\textquotesingle{} not found}

\CommentTok{\# Čitamo: R ne može naći objekt \textquotesingle{}kampanja\_tip\textquotesingle{}}
\CommentTok{\# Rješenje: provjerimo imena stupaca}
\FunctionTok{names}\NormalTok{(nl)}
\CommentTok{\# Ah, stupac se zove \textquotesingle{}campaign\_type\textquotesingle{}, ne \textquotesingle{}kampanja\_tip\textquotesingle{}}
\end{Highlighting}
\end{Shaded}

Poruka ``object not found'' gotovo uvijek znači jednu od tri stvari:
pogrešno ime (tipfeler), objekt još nije kreiran (izvršili ste kod izvan
redoslijeda) ili je objekt u drugom okruženju (na primjer, kreiran
unutar funkcije ali ne i izvan nje).

\subsection{Strategija: izoliraj
problem}\label{strategija-izoliraj-problem}

Kad imate dugački pipeline koji ne radi, razbijte ga na dijelove i
pokrenite svaki zasebno.

\begin{Shaded}
\begin{Highlighting}[]
\CommentTok{\# Umjesto pokretanja cijelog pipelinea odjednom:}
\CommentTok{\# nl |\textgreater{} filter(...) |\textgreater{} mutate(...) |\textgreater{} group\_by(...) |\textgreater{} summarise(...)}

\CommentTok{\# Pokrenite korak po korak:}
\NormalTok{korak1 }\OtherTok{\textless{}{-}}\NormalTok{ nl }\SpecialCharTok{|\textgreater{}} \FunctionTok{filter}\NormalTok{(campaign\_type }\SpecialCharTok{==} \StringTok{"weekly\_digest"}\NormalTok{)}
\NormalTok{korak1  }\CommentTok{\# Provjerite: izgleda li ovo kako očekujete?}
\end{Highlighting}
\end{Shaded}

\begin{verbatim}
# A tibble: 12 x 13
   campaign_id campaign_type subject_style   day_sent    send_hour subscribers
   <chr>       <chr>         <chr>           <chr>           <dbl>       <dbl>
 1 NL-002      weekly_digest hitno           petak              11       14266
 2 NL-007      weekly_digest upitni          srijeda            13       23450
 3 NL-009      weekly_digest personalizirani subota             11       12444
 4 NL-011      weekly_digest personalizirani utorak             11       23941
 5 NL-016      weekly_digest informativni    cetvrtak           15       12188
 6 NL-022      weekly_digest hitno           utorak             18       12768
 7 NL-026      weekly_digest brojke          utorak              7       12035
 8 NL-032      weekly_digest personalizirani utorak             18       18310
 9 NL-034      weekly_digest brojke          cetvrtak           20       24014
10 NL-035      weekly_digest hitno           ponedjeljak         7       19975
11 NL-037      weekly_digest personalizirani nedjelja           18       20070
12 NL-038      weekly_digest informativni    nedjelja           14       23842
# i 7 more variables: open_rate <dbl>, click_rate <dbl>,
#   unsubscribe_rate <dbl>, word_count <dbl>, n_links <dbl>, has_image <lgl>,
#   revenue <dbl>
\end{verbatim}

\begin{Shaded}
\begin{Highlighting}[]
\NormalTok{korak2 }\OtherTok{\textless{}{-}}\NormalTok{ korak1 }\SpecialCharTok{|\textgreater{}} \FunctionTok{mutate}\NormalTok{(}\AttributeTok{or\_pct =}\NormalTok{ open\_rate }\SpecialCharTok{*} \DecValTok{100}\NormalTok{)}
\NormalTok{korak2 }\SpecialCharTok{|\textgreater{}} \FunctionTok{select}\NormalTok{(campaign\_id, open\_rate, or\_pct) }\SpecialCharTok{|\textgreater{}} \FunctionTok{head}\NormalTok{(}\DecValTok{3}\NormalTok{)}
\end{Highlighting}
\end{Shaded}

\begin{verbatim}
# A tibble: 3 x 3
  campaign_id open_rate or_pct
  <chr>           <dbl>  <dbl>
1 NL-002          0.270   27.0
2 NL-007          0.288   28.8
3 NL-009          0.283   28.3
\end{verbatim}

\begin{Shaded}
\begin{Highlighting}[]
\CommentTok{\# OK, ovo radi. Idemo dalje...}

\NormalTok{korak3 }\OtherTok{\textless{}{-}}\NormalTok{ korak2 }\SpecialCharTok{|\textgreater{}}
  \FunctionTok{summarise}\NormalTok{(}
    \AttributeTok{n =} \FunctionTok{n}\NormalTok{(),}
    \AttributeTok{prosjek =} \FunctionTok{round}\NormalTok{(}\FunctionTok{mean}\NormalTok{(or\_pct), }\DecValTok{1}\NormalTok{)}
\NormalTok{  )}
\NormalTok{korak3}
\end{Highlighting}
\end{Shaded}

\begin{verbatim}
# A tibble: 1 x 2
      n prosjek
  <int>   <dbl>
1    12    28.9
\end{verbatim}

Pohranjivanjem svakog koraka u zasebni objekt, možete točno
identificirati na kojem koraku nastaje problem. Kad pronađete i
ispravite grešku, spojite korake natrag u pipeline.

\subsection{print() i glimpse() kao
dijagnostika}\label{print-i-glimpse-kao-dijagnostika}

Unutar funkcija i petlji, dodajte privremene \texttt{print()} naredbe da
vidite što se događa.

\begin{Shaded}
\begin{Highlighting}[]
\CommentTok{\# Debugging s print naredbama}
\NormalTok{analiziraj\_debug }\OtherTok{\textless{}{-}} \ControlFlowTok{function}\NormalTok{(data, tip) \{}
\NormalTok{  podaci }\OtherTok{\textless{}{-}}\NormalTok{ data }\SpecialCharTok{|\textgreater{}} \FunctionTok{filter}\NormalTok{(campaign\_type }\SpecialCharTok{==}\NormalTok{ tip)}
  \FunctionTok{cat}\NormalTok{(}\StringTok{"Tip:"}\NormalTok{, tip, }\StringTok{"| Redova:"}\NormalTok{, }\FunctionTok{nrow}\NormalTok{(podaci), }\StringTok{"}\SpecialCharTok{\textbackslash{}n}\StringTok{"}\NormalTok{)  }\CommentTok{\# Debug ispis}
  
  \ControlFlowTok{if}\NormalTok{ (}\FunctionTok{nrow}\NormalTok{(podaci) }\SpecialCharTok{==} \DecValTok{0}\NormalTok{) \{}
    \FunctionTok{cat}\NormalTok{(}\StringTok{"UPOZORENJE: nema podataka za tip"}\NormalTok{, tip, }\StringTok{"}\SpecialCharTok{\textbackslash{}n}\StringTok{"}\NormalTok{)}
    \FunctionTok{return}\NormalTok{(}\ConstantTok{NULL}\NormalTok{)}
\NormalTok{  \}}
  
\NormalTok{  podaci }\SpecialCharTok{|\textgreater{}}
    \FunctionTok{summarise}\NormalTok{(}
      \AttributeTok{tip =}\NormalTok{ tip,}
      \AttributeTok{n =} \FunctionTok{n}\NormalTok{(),}
      \AttributeTok{or\_M =} \FunctionTok{round}\NormalTok{(}\FunctionTok{mean}\NormalTok{(open\_rate), }\DecValTok{3}\NormalTok{)}
\NormalTok{    )}
\NormalTok{\}}

\CommentTok{\# Testirajte s poznatim i nepoznatim tipom}
\FunctionTok{analiziraj\_debug}\NormalTok{(nl, }\StringTok{"weekly\_digest"}\NormalTok{)}
\end{Highlighting}
\end{Shaded}

\begin{verbatim}
Tip: weekly_digest | Redova: 12 
\end{verbatim}

\begin{verbatim}
# A tibble: 1 x 3
  tip               n  or_M
  <chr>         <int> <dbl>
1 weekly_digest    12 0.289
\end{verbatim}

\begin{Shaded}
\begin{Highlighting}[]
\FunctionTok{analiziraj\_debug}\NormalTok{(nl, }\StringTok{"nepostojeci\_tip"}\NormalTok{)}
\end{Highlighting}
\end{Shaded}

\begin{verbatim}
Tip: nepostojeci_tip | Redova: 0 
UPOZORENJE: nema podataka za tip nepostojeci_tip 
\end{verbatim}

\begin{verbatim}
NULL
\end{verbatim}

Kad ste riješili problem, uklonite debug ispise. Ostavljanje privremenih
\texttt{cat()} i \texttt{print()} naredbi u gotovom kodu je loša praksa
jer zatrpava konzolu nepotrebnim ispisom.

\subsection{Česte greške i
rješenja}\label{ux10deste-greux161ke-i-rjeux161enja}

Pogledajmo najčešće greške koje ćete susresti i kako ih riješiti.

\begin{Shaded}
\begin{Highlighting}[]
\CommentTok{\# 1. "could not find function" {-}\textgreater{} paket nije učitan}
\FunctionTok{summarise}\NormalTok{(nl, }\AttributeTok{n =} \FunctionTok{n}\NormalTok{())}
\CommentTok{\# Rješenje: library(tidyverse) na početku}

\CommentTok{\# 2. "unexpected symbol" {-}\textgreater{} nedostaje zarez, operator ili zagrada}
\NormalTok{nl }\SpecialCharTok{|\textgreater{}}
  \FunctionTok{mutate}\NormalTok{(}\AttributeTok{x =}\NormalTok{ open\_rate }\AttributeTok{y =}\NormalTok{ click\_rate)  }\CommentTok{\# Nedostaje zarez}
\CommentTok{\# Rješenje: mutate(x = open\_rate, y = click\_rate)}

\CommentTok{\# 3. "+ ggplot" umjesto "|\textgreater{} ggplot"}
\NormalTok{nl }\SpecialCharTok{|\textgreater{}}
  \FunctionTok{filter}\NormalTok{(open\_rate }\SpecialCharTok{\textgreater{}} \FloatTok{0.2}\NormalTok{) }\SpecialCharTok{+}  \CommentTok{\# Krivo: + umjesto |\textgreater{}}
  \FunctionTok{ggplot}\NormalTok{(}\FunctionTok{aes}\NormalTok{(}\AttributeTok{x =}\NormalTok{ open\_rate))}
\CommentTok{\# Rješenje: koristiti |\textgreater{} do ggplot(), pa + za slojeve}

\CommentTok{\# 4. "object of type \textquotesingle{}closure\textquotesingle{} is not subsettable"}
\NormalTok{mean[}\DecValTok{1}\NormalTok{]  }\CommentTok{\# mean je funkcija, ne vektor}
\CommentTok{\# Rješenje: provjerite jeste li slučajno prepisali ime varijable imenom funkcije}
\end{Highlighting}
\end{Shaded}

\begin{quote}
Svaki iskusni programer bio je početnik koji je satima tražio zarez koji
nedostaje. Debugging nije znak neznanja, nego sastavni dio posla.
Razlika između početnika i iskusnog korisnika nije u tome što iskusni ne
griješe, nego u tome da imaju sustavan pristup traženju grešaka.
\end{quote}

\begin{center}\rule{0.5\linewidth}{0.5pt}\end{center}

\section{Quarto: integracija koda, teksta i rezultata}\label{sec-quarto}

Do sada ste pisali R kod u skriptama (.R datoteke) koje proizvode
tablice i grafove u konzoli. Quarto dokumenti (.qmd datoteke) omogućuju
nešto moćnije: integraciju teksta, koda i rezultata u jedan dokument
koji se renderira u HTML, PDF ili Word.

Zapravo, svako predavanje na ovom kolegiju je Quarto dokument. Tekst
koji čitate, kod koji vidite i grafovi koji se prikazuju nastaju iz
jedne .qmd datoteke.

\subsection{Struktura Quarto
dokumenta}\label{struktura-quarto-dokumenta}

\begin{Shaded}
\begin{Highlighting}[]
\CommentTok{\# Quarto dokument ima tri dijela:}

\CommentTok{\# 1. YAML zaglavlje (između {-}{-}{-} oznaka)}
\CommentTok{\# {-}{-}{-}}
\CommentTok{\# title: "Analiza newsletter kampanja"}
\CommentTok{\# author: "Ime Prezime"}
\CommentTok{\# date: today}
\CommentTok{\# format: html}
\CommentTok{\# {-}{-}{-}}

\CommentTok{\# 2. Tekst u Markdown formatu}
\CommentTok{\# \#\# Uvod}
\CommentTok{\# Ova analiza ispituje performanse naših newsletter kampanja...}

\CommentTok{\# 3. R code chunkovi (između \textasciigrave{}\textasciigrave{}\textasciigrave{} oznaka)}
\CommentTok{\# \textasciigrave{}\textasciigrave{}\textasciigrave{}\{r\}}
\CommentTok{\# library(tidyverse)}
\CommentTok{\# nl \textless{}{-} read\_csv("newsletter\_campaigns.csv")}
\CommentTok{\# \textasciigrave{}\textasciigrave{}\textasciigrave{}}
\end{Highlighting}
\end{Shaded}

Kad pokrenete \texttt{quarto\ render}, Quarto izvršava R kod, hvata
rezultate (tablice, grafove, ispis) i umeće ih u dokument zajedno s
tekstom. Rezultat je profesionalan izvještaj u kojem su analiza i
prezentacija neodvojivi.

\subsection{Chunk opcije za kontrolu
ispisa}\label{chunk-opcije-za-kontrolu-ispisa}

Opcije unutar code chunkova kontroliraju što se prikazuje u dokumentu.

\begin{Shaded}
\begin{Highlighting}[]
\CommentTok{\# echo: true   {-}\textgreater{} prikaži kod u dokumentu}
\CommentTok{\# echo: false  {-}\textgreater{} sakrij kod, prikaži samo rezultat}
\CommentTok{\# eval: true   {-}\textgreater{} izvrši kod}
\CommentTok{\# eval: false  {-}\textgreater{} ne izvršavaj (samo prikaži kod)}
\CommentTok{\# message: false {-}\textgreater{} sakrij poruke paketa}
\CommentTok{\# warning: false {-}\textgreater{} sakrij upozorenja}
\CommentTok{\# fig{-}width: 8  {-}\textgreater{} širina grafa u inčima}
\CommentTok{\# fig{-}height: 5 {-}\textgreater{} visina grafa u inčima}

\CommentTok{\# Za izvještaj klijentu: echo: false (ne želi vidjeti kod)}
\CommentTok{\# Za kolegicu analitičarku: echo: true (želi vidjeti kako ste to napravili)}
\end{Highlighting}
\end{Shaded}

Ova fleksibilnost je ključna. Isti Quarto dokument možete renderirati s
\texttt{echo:\ true} za interni tim (koji želi vidjeti kod) i s
\texttt{echo:\ false} za klijenta (koji želi samo rezultate). Mijenjate
jednu opciju u YAML zaglavlju i dobivate potpuno drugačiji dokument.

\subsection{Inline R kod}\label{inline-r-kod}

Osim code chunkova, R vrijednosti možete umetnuti direktno u tekst.

\begin{Shaded}
\begin{Highlighting}[]
\NormalTok{n\_kampanja }\OtherTok{\textless{}{-}} \FunctionTok{nrow}\NormalTok{(nl)}
\NormalTok{prosjek\_or }\OtherTok{\textless{}{-}} \FunctionTok{round}\NormalTok{(}\FunctionTok{mean}\NormalTok{(nl}\SpecialCharTok{$}\NormalTok{open\_rate) }\SpecialCharTok{*} \DecValTok{100}\NormalTok{, }\DecValTok{1}\NormalTok{)}
\NormalTok{najbolji\_tip }\OtherTok{\textless{}{-}}\NormalTok{ nl }\SpecialCharTok{|\textgreater{}}
  \FunctionTok{group\_by}\NormalTok{(campaign\_type) }\SpecialCharTok{|\textgreater{}}
  \FunctionTok{summarise}\NormalTok{(}\AttributeTok{or =} \FunctionTok{mean}\NormalTok{(open\_rate), }\AttributeTok{.groups =} \StringTok{"drop"}\NormalTok{) }\SpecialCharTok{|\textgreater{}}
  \FunctionTok{slice\_max}\NormalTok{(or) }\SpecialCharTok{|\textgreater{}}
  \FunctionTok{pull}\NormalTok{(campaign\_type)}
\end{Highlighting}
\end{Shaded}

U Quarto dokumentu biste napisali tekst poput: ``Analizirali smo 50
kampanja. Prosječni open rate je 25.5\%. Najbolji rezultat ima tip
weekly\_digest.''

Kad se dokument renderira, R vrijednosti se automatski umeću u tekst.
Ako se podaci promijene, tekst se automatski ažurira. Nikad više ne
morate ručno ažurirati brojke u izvještaju.

\subsection{Quarto vs R skripta: kad koristiti
što}\label{quarto-vs-r-skripta-kad-koristiti-ux161to}

R skripta (.R) je pravi izbor kad je cilj izračun, transformacija ili
generiranje outputa (tablice, grafovi, datoteke). Skripta je brza za
izvršavanje i laka za debugging.

Quarto dokument (.qmd) je pravi izbor kad je cilj komunikacija
rezultata. Izvještaj za klijenta, akademski rad, interna prezentacija,
kolegijalni materijal. Quarto integrira narativ i rezultate u jedinstven
dokument.

U praksi, mnogi analitičari koriste oboje: skripta za teški posao
(čišćenje, modeliranje), Quarto za prezentaciju rezultata. Skripta
generira čiste podatke i grafove, Quarto ih ugrađuje u priču.

\begin{center}\rule{0.5\linewidth}{0.5pt}\end{center}

\section{Funkcionali za složenije radne tokove}\label{sec-funkcionali}

Vratimo se purrr paketu i pogledajmo naprednije obrasce koji su korisni
u praksi.

\subsection{walk(): map() bez povratne
vrijednosti}\label{walk-map-bez-povratne-vrijednosti}

Ponekad želite izvršiti nešto za svaki element (na primjer, spremiti
graf) ali ne trebate povratnu vrijednost. \texttt{walk()} je varijanta
\texttt{map()} koja izvršava funkciju ali tiho odbacuje rezultat.

\begin{Shaded}
\begin{Highlighting}[]
\CommentTok{\# Spremi zaseban graf za svaki tip kampanje}
\NormalTok{tipovi }\OtherTok{\textless{}{-}} \FunctionTok{unique}\NormalTok{(nl}\SpecialCharTok{$}\NormalTok{campaign\_type)}

\FunctionTok{walk}\NormalTok{(tipovi, \textbackslash{}(tip) \{}
\NormalTok{  p }\OtherTok{\textless{}{-}}\NormalTok{ nl }\SpecialCharTok{|\textgreater{}}
    \FunctionTok{filter}\NormalTok{(campaign\_type }\SpecialCharTok{==}\NormalTok{ tip) }\SpecialCharTok{|\textgreater{}}
    \FunctionTok{ggplot}\NormalTok{(}\FunctionTok{aes}\NormalTok{(}\AttributeTok{x =}\NormalTok{ open\_rate)) }\SpecialCharTok{+}
    \FunctionTok{geom\_histogram}\NormalTok{(}\AttributeTok{fill =} \StringTok{"steelblue"}\NormalTok{, }\AttributeTok{color =} \StringTok{"white"}\NormalTok{, }\AttributeTok{bins =} \DecValTok{8}\NormalTok{) }\SpecialCharTok{+}
    \FunctionTok{labs}\NormalTok{(}\AttributeTok{title =} \FunctionTok{paste}\NormalTok{(}\StringTok{"Open rate:"}\NormalTok{, tip)) }\SpecialCharTok{+}
    \FunctionTok{theme\_minimal}\NormalTok{()}
  
  \FunctionTok{ggsave}\NormalTok{(}\FunctionTok{paste0}\NormalTok{(}\StringTok{"graf\_"}\NormalTok{, tip, }\StringTok{".png"}\NormalTok{), p, }\AttributeTok{width =} \DecValTok{7}\NormalTok{, }\AttributeTok{height =} \DecValTok{4}\NormalTok{)}
\NormalTok{\})}
\end{Highlighting}
\end{Shaded}

\texttt{walk()} je idiomatski R način za petlje koje proizvode popratne
efekte (side effects) poput spremanja datoteka, ispisa na konzolu ili
slanja emailova. Za razliku od \texttt{map()}, ne zatrpava konzolu
listom NULL vrijednosti.

\subsection{map2(): paralelna iteracija preko dva
vektora}\label{map2-paralelna-iteracija-preko-dva-vektora}

\begin{Shaded}
\begin{Highlighting}[]
\CommentTok{\# Dva vektora: metrike i njihovi naslovi}
\NormalTok{metrike }\OtherTok{\textless{}{-}} \FunctionTok{c}\NormalTok{(}\StringTok{"open\_rate"}\NormalTok{, }\StringTok{"click\_rate"}\NormalTok{)}
\NormalTok{naslovi }\OtherTok{\textless{}{-}} \FunctionTok{c}\NormalTok{(}\StringTok{"Open rate kampanja"}\NormalTok{, }\StringTok{"Click rate kampanja"}\NormalTok{)}

\CommentTok{\# map2 iterira paralelno: prvi element s prvim, drugi s drugim}
\NormalTok{rezultati }\OtherTok{\textless{}{-}} \FunctionTok{map2}\NormalTok{(metrike, naslovi, \textbackslash{}(metrika, naslov) \{}
\NormalTok{  nl }\SpecialCharTok{|\textgreater{}}
    \FunctionTok{sazetak\_metrike}\NormalTok{(metrika) }\SpecialCharTok{|\textgreater{}}
    \FunctionTok{mutate}\NormalTok{(}\AttributeTok{metrika =}\NormalTok{ naslov)}
\NormalTok{\})}

\FunctionTok{bind\_rows}\NormalTok{(rezultati)}
\end{Highlighting}
\end{Shaded}

\begin{verbatim}
# A tibble: 2 x 6
      M   Med    SD   Min   Max metrika            
  <dbl> <dbl> <dbl> <dbl> <dbl> <chr>              
1 0.255 0.252 0.068 0.115 0.408 Open rate kampanja 
2 0.037 0.037 0.019 0.005 0.086 Click rate kampanja
\end{verbatim}

\texttt{map2()} prima dva vektora i iterira paralelno. Na prvoj
iteraciji koristi \texttt{metrike{[}1{]}} i \texttt{naslovi{[}1{]}}, na
drugoj \texttt{metrike{[}2{]}} i \texttt{naslovi{[}2{]}}, i tako dalje.
Ovo je korisno kad imate parove ulaznih podataka.

\subsection{imap(): iteracija s
indeksom}\label{imap-iteracija-s-indeksom}

\begin{Shaded}
\begin{Highlighting}[]
\CommentTok{\# imap daje i element i njegovo ime/indeks}
\NormalTok{nl }\SpecialCharTok{|\textgreater{}}
  \FunctionTok{group\_by}\NormalTok{(campaign\_type) }\SpecialCharTok{|\textgreater{}}
  \FunctionTok{group\_split}\NormalTok{() }\SpecialCharTok{|\textgreater{}}
  \FunctionTok{set\_names}\NormalTok{(}\FunctionTok{unique}\NormalTok{(nl}\SpecialCharTok{$}\NormalTok{campaign\_type) }\SpecialCharTok{|\textgreater{}} \FunctionTok{sort}\NormalTok{()) }\SpecialCharTok{|\textgreater{}}
  \FunctionTok{imap}\NormalTok{(\textbackslash{}(podaci, ime) \{}
    \FunctionTok{tibble}\NormalTok{(}
      \AttributeTok{tip =}\NormalTok{ ime,}
      \AttributeTok{n =} \FunctionTok{nrow}\NormalTok{(podaci),}
      \AttributeTok{or\_M =} \FunctionTok{round}\NormalTok{(}\FunctionTok{mean}\NormalTok{(podaci}\SpecialCharTok{$}\NormalTok{open\_rate), }\DecValTok{3}\NormalTok{)}
\NormalTok{    )}
\NormalTok{  \}) }\SpecialCharTok{|\textgreater{}}
  \FunctionTok{bind\_rows}\NormalTok{()}
\end{Highlighting}
\end{Shaded}

\begin{verbatim}
# A tibble: 5 x 3
  tip                n  or_M
  <chr>          <int> <dbl>
1 breaking_news      7 0.198
2 event_promo        6 0.25 
3 special_report    17 0.259
4 sponsored          8 0.248
5 weekly_digest     12 0.289
\end{verbatim}

\texttt{imap()} je varijanta \texttt{map()} koja automatski proslijeđuje
i element i njegovo ime (ili indeks). Korisna je kad trebate znati koji
element trenutno obrađujete, na primjer za imenovanje rezultata ili za
dijagnostiku.

\subsection{possibly(): zaštita od
grešaka}\label{possibly-zaux161tita-od-greux161aka}

Kad primjenjujete funkciju na mnogo elemenata, jedna greška može srušiti
cijeli pipeline. \texttt{possibly()} omotava funkciju u zaštitni sloj
koji hvata greške i vraća default vrijednost umjesto da prekida
izvršavanje.

\begin{Shaded}
\begin{Highlighting}[]
\CommentTok{\# Funkcija koja ponekad pada}
\NormalTok{opasna\_funkcija }\OtherTok{\textless{}{-}} \ControlFlowTok{function}\NormalTok{(tip) \{}
\NormalTok{  podaci }\OtherTok{\textless{}{-}}\NormalTok{ nl }\SpecialCharTok{|\textgreater{}} \FunctionTok{filter}\NormalTok{(campaign\_type }\SpecialCharTok{==}\NormalTok{ tip)}
  \ControlFlowTok{if}\NormalTok{ (}\FunctionTok{nrow}\NormalTok{(podaci) }\SpecialCharTok{\textless{}} \DecValTok{3}\NormalTok{) }\FunctionTok{stop}\NormalTok{(}\StringTok{"Premalo podataka!"}\NormalTok{)}
  \FunctionTok{mean}\NormalTok{(podaci}\SpecialCharTok{$}\NormalTok{open\_rate)}
\NormalTok{\}}

\CommentTok{\# Bez zaštite: jedna greška ruši sve}
\CommentTok{\# map\_dbl(c("weekly\_digest", "nepostojeci"), opasna\_funkcija)  \# Error!}

\CommentTok{\# S zaštitom: greška vraća NA, ostali rezultati ostaju}
\NormalTok{sigurna\_funkcija }\OtherTok{\textless{}{-}} \FunctionTok{possibly}\NormalTok{(opasna\_funkcija, }\AttributeTok{otherwise =} \ConstantTok{NA\_real\_}\NormalTok{)}

\FunctionTok{map\_dbl}\NormalTok{(}\FunctionTok{c}\NormalTok{(}\StringTok{"weekly\_digest"}\NormalTok{, }\StringTok{"nepostojeci"}\NormalTok{, }\StringTok{"breaking\_news"}\NormalTok{), sigurna\_funkcija)}
\end{Highlighting}
\end{Shaded}

\begin{verbatim}
[1] 0.2886833        NA 0.1978286
\end{verbatim}

\texttt{possibly(f,\ otherwise\ =\ NA)} kreira novu funkciju koja radi
isto kao \texttt{f}, ali umjesto da baci grešku, vraća
\texttt{otherwise} vrijednost. Ovo je neprocjenjivo kad učitavate 50
datoteka i jedna je korumpirana, ili kad analizirate 20 grupa i jedna
ima nedovoljno podataka.

\begin{center}\rule{0.5\linewidth}{0.5pt}\end{center}

\section{Kompletna analiza: automatizirani izvještaj o
kampanjama}\label{sec-kompletna}

Spojimo sve iz ovog predavanja u jednu koherentnu analizu. Cilj je
napisati kod koji bi mogao biti tijelo Quarto izvještaja o performansama
newsletter kampanja.

\begin{Shaded}
\begin{Highlighting}[]
\FunctionTok{library}\NormalTok{(patchwork)}

\CommentTok{\# PARAMETRI}
\NormalTok{min\_kampanja }\OtherTok{\textless{}{-}} \DecValTok{3}
\NormalTok{decimale }\OtherTok{\textless{}{-}} \DecValTok{3}
\NormalTok{fokus\_metrike }\OtherTok{\textless{}{-}} \FunctionTok{c}\NormalTok{(}\StringTok{"open\_rate"}\NormalTok{, }\StringTok{"click\_rate"}\NormalTok{, }\StringTok{"unsubscribe\_rate"}\NormalTok{)}
\end{Highlighting}
\end{Shaded}

\begin{Shaded}
\begin{Highlighting}[]
\CommentTok{\# POMOĆNE FUNKCIJE}
\NormalTok{sazetak\_tipa }\OtherTok{\textless{}{-}} \ControlFlowTok{function}\NormalTok{(data, tip, }\AttributeTok{dec =} \DecValTok{3}\NormalTok{) \{}
\NormalTok{  d }\OtherTok{\textless{}{-}}\NormalTok{ data }\SpecialCharTok{|\textgreater{}} \FunctionTok{filter}\NormalTok{(campaign\_type }\SpecialCharTok{==}\NormalTok{ tip)}
  
  \ControlFlowTok{if}\NormalTok{ (}\FunctionTok{nrow}\NormalTok{(d) }\SpecialCharTok{\textless{}}\NormalTok{ min\_kampanja) }\FunctionTok{return}\NormalTok{(}\ConstantTok{NULL}\NormalTok{)}
  
  \FunctionTok{tibble}\NormalTok{(}
    \AttributeTok{tip =}\NormalTok{ tip,}
    \AttributeTok{n =} \FunctionTok{nrow}\NormalTok{(d),}
    \AttributeTok{or\_M =} \FunctionTok{round}\NormalTok{(}\FunctionTok{mean}\NormalTok{(d}\SpecialCharTok{$}\NormalTok{open\_rate), dec),}
    \AttributeTok{or\_SD =} \FunctionTok{round}\NormalTok{(}\FunctionTok{sd}\NormalTok{(d}\SpecialCharTok{$}\NormalTok{open\_rate), dec),}
    \AttributeTok{ctr\_M =} \FunctionTok{round}\NormalTok{(}\FunctionTok{mean}\NormalTok{(d}\SpecialCharTok{$}\NormalTok{click\_rate), dec }\SpecialCharTok{+} \DecValTok{1}\NormalTok{),}
    \AttributeTok{unsub\_M =} \FunctionTok{round}\NormalTok{(}\FunctionTok{mean}\NormalTok{(d}\SpecialCharTok{$}\NormalTok{unsubscribe\_rate), dec }\SpecialCharTok{+} \DecValTok{2}\NormalTok{),}
    \AttributeTok{prosj\_rijeci =} \FunctionTok{round}\NormalTok{(}\FunctionTok{mean}\NormalTok{(d}\SpecialCharTok{$}\NormalTok{word\_count), }\DecValTok{0}\NormalTok{),}
    \AttributeTok{udio\_sa\_slikom =} \FunctionTok{round}\NormalTok{(}\FunctionTok{mean}\NormalTok{(d}\SpecialCharTok{$}\NormalTok{has\_image), }\DecValTok{2}\NormalTok{)}
\NormalTok{  )}
\NormalTok{\}}

\NormalTok{graf\_usporedba }\OtherTok{\textless{}{-}} \ControlFlowTok{function}\NormalTok{(data, metrika, naslov, }\AttributeTok{boja =} \StringTok{"steelblue"}\NormalTok{) \{}
\NormalTok{  data }\SpecialCharTok{|\textgreater{}}
    \FunctionTok{ggplot}\NormalTok{(}\FunctionTok{aes}\NormalTok{(}\AttributeTok{x =} \FunctionTok{fct\_reorder}\NormalTok{(campaign\_type, .data[[metrika]]),}
               \AttributeTok{y =}\NormalTok{ .data[[metrika]])) }\SpecialCharTok{+}
    \FunctionTok{geom\_boxplot}\NormalTok{(}\AttributeTok{fill =}\NormalTok{ boja, }\AttributeTok{alpha =} \FloatTok{0.6}\NormalTok{) }\SpecialCharTok{+}
    \FunctionTok{coord\_flip}\NormalTok{() }\SpecialCharTok{+}
    \FunctionTok{labs}\NormalTok{(}\AttributeTok{title =}\NormalTok{ naslov, }\AttributeTok{x =} \ConstantTok{NULL}\NormalTok{, }\AttributeTok{y =}\NormalTok{ metrika) }\SpecialCharTok{+}
    \FunctionTok{theme\_minimal}\NormalTok{()}
\NormalTok{\}}
\end{Highlighting}
\end{Shaded}

\begin{Shaded}
\begin{Highlighting}[]
\CommentTok{\# ANALIZA}
\NormalTok{tipovi }\OtherTok{\textless{}{-}} \FunctionTok{unique}\NormalTok{(nl}\SpecialCharTok{$}\NormalTok{campaign\_type)}

\CommentTok{\# Sažetak za sve tipove (s automatskim preskakanjem malih grupa)}
\NormalTok{tablica\_sazetka }\OtherTok{\textless{}{-}} \FunctionTok{map}\NormalTok{(tipovi, \textbackslash{}(t) }\FunctionTok{sazetak\_tipa}\NormalTok{(nl, t, decimale)) }\SpecialCharTok{|\textgreater{}}
  \FunctionTok{bind\_rows}\NormalTok{() }\SpecialCharTok{|\textgreater{}}
  \FunctionTok{arrange}\NormalTok{(}\FunctionTok{desc}\NormalTok{(or\_M))}

\NormalTok{tablica\_sazetka}
\end{Highlighting}
\end{Shaded}

\begin{verbatim}
# A tibble: 5 x 8
  tip                n  or_M or_SD  ctr_M unsub_M prosj_rijeci udio_sa_slikom
  <chr>          <int> <dbl> <dbl>  <dbl>   <dbl>        <dbl>          <dbl>
1 weekly_digest     12 0.289 0.066 0.029  0.00331          380           0.67
2 special_report    17 0.259 0.067 0.0519 0.00279          486           0.76
3 event_promo        6 0.25  0.086 0.0294 0.00207          210           1   
4 sponsored          8 0.248 0.043 0.0223 0.00645          204           0.75
5 breaking_news      7 0.198 0.055 0.0382 0.00405          140           0.71
\end{verbatim}

\begin{Shaded}
\begin{Highlighting}[]
\CommentTok{\# Analiza po stilu naslova (unutar svake kampanje)}
\NormalTok{nl }\SpecialCharTok{|\textgreater{}}
  \FunctionTok{group\_by}\NormalTok{(campaign\_type, subject\_style) }\SpecialCharTok{|\textgreater{}}
  \FunctionTok{summarise}\NormalTok{(}
    \AttributeTok{n =} \FunctionTok{n}\NormalTok{(),}
    \AttributeTok{or\_M =} \FunctionTok{round}\NormalTok{(}\FunctionTok{mean}\NormalTok{(open\_rate), }\DecValTok{3}\NormalTok{),}
    \AttributeTok{.groups =} \StringTok{"drop"}
\NormalTok{  ) }\SpecialCharTok{|\textgreater{}}
  \FunctionTok{filter}\NormalTok{(n }\SpecialCharTok{\textgreater{}=} \DecValTok{2}\NormalTok{) }\SpecialCharTok{|\textgreater{}}
  \FunctionTok{pivot\_wider}\NormalTok{(}
    \AttributeTok{names\_from =}\NormalTok{ subject\_style,}
    \AttributeTok{values\_from =}\NormalTok{ or\_M}
\NormalTok{  )}
\end{Highlighting}
\end{Shaded}

\begin{verbatim}
# A tibble: 12 x 7
   campaign_type      n informativni personalizirani upitni brojke  hitno
   <chr>          <int>        <dbl>           <dbl>  <dbl>  <dbl>  <dbl>
 1 breaking_news      5        0.182          NA     NA     NA     NA    
 2 event_promo        2       NA               0.285 NA     NA     NA    
 3 event_promo        3       NA              NA      0.245 NA     NA    
 4 special_report     2       NA               0.242 NA      0.237 NA    
 5 special_report     4       NA              NA     NA     NA      0.329
 6 special_report     6        0.2            NA     NA     NA     NA    
 7 special_report     3       NA              NA      0.31  NA     NA    
 8 sponsored          3       NA              NA     NA      0.252 NA    
 9 sponsored          2       NA              NA     NA     NA      0.224
10 weekly_digest      2        0.21           NA     NA      0.254 NA    
11 weekly_digest      3       NA              NA     NA     NA      0.312
12 weekly_digest      4       NA               0.328 NA     NA     NA    
\end{verbatim}

\begin{Shaded}
\begin{Highlighting}[]
\CommentTok{\# VIZUALIZACIJA}
\NormalTok{g1 }\OtherTok{\textless{}{-}} \FunctionTok{graf\_usporedba}\NormalTok{(nl, }\StringTok{"open\_rate"}\NormalTok{, }\StringTok{"Open rate po tipu"}\NormalTok{)}
\NormalTok{g2 }\OtherTok{\textless{}{-}} \FunctionTok{graf\_usporedba}\NormalTok{(nl, }\StringTok{"click\_rate"}\NormalTok{, }\StringTok{"Click rate po tipu"}\NormalTok{, }\AttributeTok{boja =} \StringTok{"\#2a9d8f"}\NormalTok{)}

\NormalTok{g1 }\SpecialCharTok{+}\NormalTok{ g2 }\SpecialCharTok{+}
  \FunctionTok{plot\_annotation}\NormalTok{(}
    \AttributeTok{title =} \StringTok{"Performanse newsletter kampanja"}\NormalTok{,}
    \AttributeTok{subtitle =} \StringTok{"Usporedba open rate i click rate po tipu kampanje"}\NormalTok{,}
    \AttributeTok{caption =} \FunctionTok{paste}\NormalTok{(}\StringTok{"N ="}\NormalTok{, }\FunctionTok{nrow}\NormalTok{(nl), }\StringTok{"kampanja"}\NormalTok{)}
\NormalTok{  )}
\end{Highlighting}
\end{Shaded}

\pandocbounded{\includegraphics[keepaspectratio]{week-06_files/figure-pdf/kompletna-vizualizacija-1.pdf}}

\begin{Shaded}
\begin{Highlighting}[]
\CommentTok{\# Odnos open rate i click rate}
\NormalTok{nl }\SpecialCharTok{|\textgreater{}}
  \FunctionTok{ggplot}\NormalTok{(}\FunctionTok{aes}\NormalTok{(}\AttributeTok{x =}\NormalTok{ open\_rate, }\AttributeTok{y =}\NormalTok{ click\_rate, }\AttributeTok{color =}\NormalTok{ campaign\_type, }\AttributeTok{size =}\NormalTok{ subscribers)) }\SpecialCharTok{+}
  \FunctionTok{geom\_point}\NormalTok{(}\AttributeTok{alpha =} \FloatTok{0.6}\NormalTok{) }\SpecialCharTok{+}
  \FunctionTok{geom\_smooth}\NormalTok{(}\AttributeTok{method =} \StringTok{"lm"}\NormalTok{, }\AttributeTok{se =} \ConstantTok{FALSE}\NormalTok{, }\AttributeTok{linewidth =} \FloatTok{0.8}\NormalTok{) }\SpecialCharTok{+}
  \FunctionTok{scale\_size\_continuous}\NormalTok{(}\AttributeTok{labels =}\NormalTok{ scales}\SpecialCharTok{::}\FunctionTok{label\_comma}\NormalTok{()) }\SpecialCharTok{+}
  \FunctionTok{labs}\NormalTok{(}
    \AttributeTok{title =} \StringTok{"Kampanje s višim open rateom tendiraju imati viši click rate"}\NormalTok{,}
    \AttributeTok{subtitle =} \StringTok{"Veličina točke proporcionalna broju pretplatnika"}\NormalTok{,}
    \AttributeTok{x =} \StringTok{"Open rate"}\NormalTok{,}
    \AttributeTok{y =} \StringTok{"Click rate"}\NormalTok{,}
    \AttributeTok{color =} \StringTok{"Tip kampanje"}\NormalTok{,}
    \AttributeTok{size =} \StringTok{"Pretplatnici"}
\NormalTok{  ) }\SpecialCharTok{+}
  \FunctionTok{theme\_minimal}\NormalTok{() }\SpecialCharTok{+}
  \FunctionTok{theme}\NormalTok{(}\AttributeTok{legend.position =} \StringTok{"bottom"}\NormalTok{)}
\end{Highlighting}
\end{Shaded}

\pandocbounded{\includegraphics[keepaspectratio]{week-06_files/figure-pdf/kompletna-scatter-1.pdf}}

\begin{Shaded}
\begin{Highlighting}[]
\CommentTok{\# Kada slati newsletter?}
\NormalTok{nl }\SpecialCharTok{|\textgreater{}}
  \FunctionTok{mutate}\NormalTok{(}
    \AttributeTok{dio\_dana =} \FunctionTok{case\_when}\NormalTok{(}
\NormalTok{      send\_hour }\SpecialCharTok{\textless{}} \DecValTok{10} \SpecialCharTok{\textasciitilde{}} \StringTok{"jutro (6{-}9)"}\NormalTok{,}
\NormalTok{      send\_hour }\SpecialCharTok{\textless{}} \DecValTok{14} \SpecialCharTok{\textasciitilde{}} \StringTok{"prijepodne (10{-}13)"}\NormalTok{,}
\NormalTok{      send\_hour }\SpecialCharTok{\textless{}} \DecValTok{18} \SpecialCharTok{\textasciitilde{}} \StringTok{"poslijepodne (14{-}17)"}\NormalTok{,}
      \AttributeTok{.default =} \StringTok{"navečer (18+)"}
\NormalTok{    )}
\NormalTok{  ) }\SpecialCharTok{|\textgreater{}}
  \FunctionTok{group\_by}\NormalTok{(dio\_dana) }\SpecialCharTok{|\textgreater{}}
  \FunctionTok{summarise}\NormalTok{(}
    \AttributeTok{n =} \FunctionTok{n}\NormalTok{(),}
    \AttributeTok{or\_M =} \FunctionTok{round}\NormalTok{(}\FunctionTok{mean}\NormalTok{(open\_rate), }\DecValTok{3}\NormalTok{),}
    \AttributeTok{ctr\_M =} \FunctionTok{round}\NormalTok{(}\FunctionTok{mean}\NormalTok{(click\_rate), }\DecValTok{4}\NormalTok{),}
    \AttributeTok{.groups =} \StringTok{"drop"}
\NormalTok{  ) }\SpecialCharTok{|\textgreater{}}
  \FunctionTok{arrange}\NormalTok{(}\FunctionTok{desc}\NormalTok{(or\_M))}
\end{Highlighting}
\end{Shaded}

\begin{verbatim}
# A tibble: 4 x 4
  dio_dana                 n  or_M  ctr_M
  <chr>                <int> <dbl>  <dbl>
1 navečer (18+)            9 0.277 0.0388
2 poslijepodne (14-17)    15 0.26  0.04  
3 prijepodne (10-13)      10 0.245 0.0307
4 jutro (6-9)             16 0.243 0.0373
\end{verbatim}

\begin{Shaded}
\begin{Highlighting}[]
\CommentTok{\# KLJUČNI NALAZ: koji faktori predviđaju open rate?}
\NormalTok{nl }\SpecialCharTok{|\textgreater{}}
  \FunctionTok{summarise}\NormalTok{(}
    \AttributeTok{kor\_rijeci\_or =} \FunctionTok{round}\NormalTok{(}\FunctionTok{cor}\NormalTok{(word\_count, open\_rate), }\DecValTok{3}\NormalTok{),}
    \AttributeTok{kor\_linkovi\_ctr =} \FunctionTok{round}\NormalTok{(}\FunctionTok{cor}\NormalTok{(n\_links, click\_rate), }\DecValTok{3}\NormalTok{),}
    \AttributeTok{kor\_pretplatnici\_or =} \FunctionTok{round}\NormalTok{(}\FunctionTok{cor}\NormalTok{(subscribers, open\_rate), }\DecValTok{3}\NormalTok{)}
\NormalTok{  )}
\end{Highlighting}
\end{Shaded}

\begin{verbatim}
# A tibble: 1 x 3
  kor_rijeci_or kor_linkovi_ctr kor_pretplatnici_or
          <dbl>           <dbl>               <dbl>
1         0.127          -0.239              -0.075
\end{verbatim}

Analiza otkriva nekoliko nalaza. Breaking news i kampanje s hitnim
stilom naslova imaju najviši open rate, ali uz veću varijabilnost. Click
rate ne prati savršeno open rate, što sugerira da su čimbenici koji
navode ljude da otvore email (naslov, hitnost) različiti od onih koji ih
navode da kliknu na sadržaj (relevantnost, format). Korelacija između
broja riječi i open ratea govori o preferiranom formatu, dok veličina
baze pretplatnika sama po sebi ne predviđa bolji angažman.

Cijela ova analiza, od učitavanja podataka do gotovih grafova i nalaza,
koristi principe koje smo naučili ovaj tjedan. Parametri su na jednom
mjestu. Pomoćne funkcije enkapsuliraju ponovljivu logiku. map()
automatizira iteraciju. Vizualizacija prati principe iz prošlog tjedna.
I sve je napisano tako da se može pokrenuti ponovno s novim podacima bez
ikakvih promjena u kodu (osim, eventualno, putanje do datoteke).

\begin{quote}
Kad pišete analizu, zamislite da ju netko pokreće prvi put, bez ikakvog
konteksta. Može li taj netko razumjeti što kod radi, zašto, i kako
interpretirati rezultate? Ako da, napisali ste dobru analizu.
\end{quote}

\begin{center}\rule{0.5\linewidth}{0.5pt}\end{center}

\begin{tcolorbox}[enhanced jigsaw, coltitle=black, arc=.35mm, colframe=quarto-callout-important-color-frame, left=2mm, bottomrule=.15mm, toprule=.15mm, leftrule=.75mm, bottomtitle=1mm, colback=white, title=\textcolor{quarto-callout-important-color}{\faExclamation}\hspace{0.5em}{Ključni zaključci}, opacityback=0, colbacktitle=quarto-callout-important-color!10!white, toptitle=1mm, titlerule=0mm, opacitybacktitle=0.6, breakable, rightrule=.15mm]

\begin{enumerate}
\def\labelenumi{\arabic{enumi}.}
\item
  Funkcije su alat za izbjegavanje ponavljanja koda. Pravilo tri: ako
  ste kopirali isti kod tri puta, pretvorite ga u funkciju. Default
  vrijednosti argumenata čine funkcije fleksibilnima.
\item
  Klasični \texttt{if}/\texttt{else} radi s jednom vrijednošću (za
  skripte i funkcije). \texttt{if\_else()} i \texttt{case\_when()} su
  vektorizirani (za \texttt{mutate()}). Ne miješajte ih.
\item
  Validacija ulaza u funkcijama sprečava tihe greške. Koristite
  \texttt{warning()} za upozorenja i \texttt{return()} za rano
  izlaženje.
\item
  \texttt{for} petlje ponavljaju kod za svaki element. Unutar petlje,
  grafove morate ispisati s \texttt{print()}. Za većinu zadataka postoje
  elegantnije alternative.
\item
  \texttt{map()} iz paketa purrr je moderna alternativa petljama.
  \texttt{map\_dbl()}, \texttt{map\_chr()} i \texttt{map\_lgl()} vraćaju
  specifične tipove. \texttt{walk()} je za popratne efekte (spremanje
  datoteka).
\item
  \texttt{map2()} iterira paralelno preko dva vektora. \texttt{imap()}
  daje i element i njegovo ime. \texttt{possibly()} štiti od grešaka
  unutar iteracije.
\item
  Obrazac
  \texttt{list.files()\ \textbar{}\textgreater{}\ map(read\_csv)\ \textbar{}\textgreater{}\ bind\_rows()}
  učitava i spaja više datoteka u jednom koraku.
\item
  Debugging zahtijeva sustavan pristup: čitanje poruka, izolacija
  problema (korak po korak), privremeni \texttt{cat()}/\texttt{print()}
  ispisi.
\item
  DRY princip: parametri na jednom mjestu, logika u funkcijama,
  struktura skripte u jasnim sekcijama.
\item
  Quarto dokumenti integriraju tekst, kod i rezultate. Koristite ih za
  izvještaje, radove i prezentacije. R skripte su za teški izračun,
  Quarto za komunikaciju.
\item
  Chunk opcije (\texttt{echo}, \texttt{eval}, \texttt{message},
  \texttt{warning}, \texttt{fig-width}) kontroliraju što se prikazuje u
  renderiranom dokumentu. \texttt{echo:\ false} sakriva kod za klijente.
\item
  Cilj ponovljive analize: netko može pokrenuti vaš kod od početka do
  kraja s novim podacima i dobiti ažurirane rezultate bez ručnih
  promjena.
\end{enumerate}

\end{tcolorbox}

\begin{center}\rule{0.5\linewidth}{0.5pt}\end{center}

\begin{tcolorbox}[enhanced jigsaw, coltitle=black, arc=.35mm, colframe=quarto-callout-warning-color-frame, left=2mm, bottomrule=.15mm, toprule=.15mm, leftrule=.75mm, bottomtitle=1mm, colback=white, title=\textcolor{quarto-callout-warning-color}{\faExclamationTriangle}\hspace{0.5em}{Priprema za sljedeći tjedan}, opacityback=0, colbacktitle=quarto-callout-warning-color!10!white, toptitle=1mm, titlerule=0mm, opacitybacktitle=0.6, breakable, rightrule=.15mm]

Sljedeći tjedan ulazimo u \textbf{uvod u vjerojatnost}: što je
vjerojatnost, kako ju računamo, binomna i normalna distribucija. Ovo je
konceptualni temelj za sve statističke testove koje ćemo raditi u drugom
dijelu kolegija.

Za pripremu:

\begin{enumerate}
\def\labelenumi{\arabic{enumi}.}
\tightlist
\item
  Napišite vlastitu funkciju koja prima tibble i ime kategoričke
  varijable te vraća tibble s brojem i udjelom (\%) svake kategorije.
  Testirajte je na datasetu \texttt{newsletter\_campaigns.csv}.
\item
  Koristeći \texttt{map()}, generirajte sažetak open ratea za svaki dan
  u tjednu (stupac \texttt{day\_sent}). Spojite rezultate u jedan
  tibble.
\item
  Napišite kratki Quarto dokument (.qmd) koji učitava podatke, prikazuje
  jedan graf i jednu tablicu, s popratnim tekstom. Renderirajte ga u
  HTML.
\item
  Pročitajte poglavlje 9 iz Navarro (Learning Statistics with R) o
  vjerojatnosti. Fokusirajte se na intuiciju, ne na formule.
\end{enumerate}

\end{tcolorbox}

\begin{center}\rule{0.5\linewidth}{0.5pt}\end{center}

\section{Dodatno čitanje}\label{dodatno-ux10ditanje}

\textbf{Obavezno}

Wickham, H. \& Grolemund, G. (2023). \emph{R for Data Science} (2nd
edition), Chapters 26, 27 i 29. Besplatno dostupno na
\href{https://r4ds.hadley.nz}{r4ds.hadley.nz}. Poglavlje 26 pokriva
funkcije, poglavlje 27 iteraciju s purrr, poglavlje 29 Quarto dokumente.

Navarro, D. (2018). \emph{Learning Statistics with R}, Chapter 8.
Besplatno dostupno na
\href{https://learningstatisticswithr.com/lsr-0.6.pdf}{learningstatisticswithr.com}.
Osnove programiranja u R-u.

\textbf{Preporučeno}

Wickham, H. (2019). \emph{Advanced R} (2nd edition), Chapters 6 i 9.
Besplatno dostupno na \href{https://adv-r.hadley.nz}{adv-r.hadley.nz}.
Poglavlje 6 detaljno pokriva funkcije, poglavlje 9 funkcionalno
programiranje (map i prijatelji).

Quarto dokumentacija. Besplatno dostupno na
\href{https://quarto.org}{quarto.org}. Kompletna dokumentacija za Quarto
sustav sa tutorijalima za HTML, PDF i Word dokumente.

Bryan, J. \& Hester, J. \emph{What They Forgot to Teach You About R}.
Besplatno dostupno na \href{https://rstats.wtf}{rstats.wtf}. Praktični
savjeti o organizaciji projekata, debugging-u i radnim tokovima koji se
ne uče u udžbenicima statistike.

\begin{center}\rule{0.5\linewidth}{0.5pt}\end{center}

\section{Pojmovnik}\label{pojmovnik}

\begin{longtable}[]{@{}
  >{\raggedright\arraybackslash}p{(\linewidth - 2\tabcolsep) * \real{0.5000}}
  >{\raggedright\arraybackslash}p{(\linewidth - 2\tabcolsep) * \real{0.5000}}@{}}
\toprule\noalign{}
\begin{minipage}[b]{\linewidth}\raggedright
Pojam
\end{minipage} & \begin{minipage}[b]{\linewidth}\raggedright
Objašnjenje
\end{minipage} \\
\midrule\noalign{}
\endhead
\bottomrule\noalign{}
\endlastfoot
Funkcija & Objekt koji prima argumente, izvršava operacije i vraća
rezultat. Definira se s \texttt{function()}. \\
Argument & Ulazni podatak funkcije. Navodi se unutar zagrada pri
definiciji i pozivu. \\
Default vrijednost & Podrazumijevana vrijednost argumenta. Definira se s
\texttt{=} u listi argumenata. \\
Povratna vrijednost & Rezultat funkcije. Zadnji izraz u tijelu, ili
eksplicitno s \texttt{return()}. \\
\texttt{return()} & Eksplicitno vraća vrijednost i izlazi iz funkcije.
Korisno za ranu validaciju. \\
\texttt{if}/\texttt{else} & Uvjetna naredba za kontrolu toka. Radi s
jednom vrijednošću (nije vektorizirana). \\
\texttt{if\_else()} & Vektorizirana uvjetna funkcija za
\texttt{mutate()}. Radi na cijelom stupcu. \\
\texttt{case\_when()} & Vektorizirana funkcija za složeno rekodiranje s
više uvjeta. \\
\texttt{for} petlja & Ponavlja blok koda za svaki element u skupu.
Sintaksa: \texttt{for\ (x\ in\ skup)\ \{\ ...\ \}}. \\
\texttt{map()} & purrr funkcija koja primjenjuje funkciju na svaki
element i vraća listu. \\
\texttt{map\_dbl()} & Varijanta \texttt{map()} koja vraća numerički
vektor. \\
\texttt{map\_chr()} & Varijanta \texttt{map()} koja vraća tekstualni
vektor. \\
\texttt{map\_lgl()} & Varijanta \texttt{map()} koja vraća logički
vektor. \\
\texttt{map2()} & purrr funkcija za paralelnu iteraciju preko dva
vektora. \\
\texttt{imap()} & purrr funkcija koja proslijeđuje i element i njegovo
ime/indeks. \\
\texttt{walk()} & Varijanta \texttt{map()} za popratne efekte (spremanje
datoteka). Ne vraća rezultat. \\
\texttt{possibly()} & purrr funkcija koja omotava funkciju u zaštitni
sloj. Greška vraća default vrijednost umjesto prekida. \\
purrr & Paket iz tidyverse za funkcijsko programiranje. \\
\texttt{nest()} & tidyr funkcija koja pakira podatke grupe u ugniježđeni
tibble. \\
\texttt{bind\_rows()} & Vertikalno spaja listu tibbleova u jedan. \\
\texttt{list.files()} & Base R funkcija za pronalaženje datoteka u
direktoriju po uzorku. \\
\texttt{set\_names()} & Dodjeljuje imena elementima vektora ili
liste. \\
DRY & Don't Repeat Yourself. Princip da informacija postoji na jednom
mjestu u kodu. \\
Lambda funkcija & Anonimna funkcija. Piše se kao
\texttt{\textbackslash{}(x)\ x\ +\ 1} ili
\texttt{function(x)\ x\ +\ 1}. \\
\texttt{.data{[}{[}var{]}{]}} & Pristup stupcu po imenu pohranjenom u
varijabli. Za tidyverse funkcije. \\
\texttt{cat()} & Ispis teksta u konzolu. Bez navodnih oznaka i
indeksa. \\
\texttt{warning()} & Ispis upozorenja. Ne zaustavlja program. \\
\texttt{stop()} & Ispis greške i zaustavljanje programa. Za kritične
probleme. \\
Validacija ulaza & Provjera ispravnosti argumenata prije izvršavanja.
Sprečava tihe greške. \\
Skripta (.R) & R datoteka s nizom naredbi. Za izračune i
transformacije. \\
Quarto dokument (.qmd) & Datoteka koja integrira tekst, kod i rezultate.
Za izvještaje i komunikaciju. \\
Chunk opcije & Postavke R code chunka u Quarto dokumentu (echo, eval,
message, warning, fig-width). \\
Inline R kod & R izraz umetnut u tekst Quarto dokumenta. Automatski se
evaluira pri renderiranju. \\
Debugging & Proces pronalaženja i ispravljanja grešaka u kodu. \\
Side effect & Popratni efekt funkcije (ispis, spremanje datoteke) koji
nije povratna vrijednost. \\
Ponovljiva analiza & Analiza napisana tako da se može pokrenuti od
početka do kraja s novim podacima bez ručnih promjena. \\
\end{longtable}




\end{document}
